\section{Часть I. Элементарная теория погрешностей}

\subsection{Источники и классификация погрешностей. Абсолютная и относительная погрешности. Значащие цифры в десятичной записи числа. Округление. Верные значащие цифры. Правила записи приближённых чисел}

\subsubsection*{Источники и классификация погрешностей}

Источники погрешностей:
\begin{enumerate}
    \item Приближённость математической модели (погрешность модели).
    \item Погрешности исходных данных (измерений).
    \item Приближённость метода численного решения (погрешность метода).
    \item Погрешности машинного представления чисел и арифметических операций (вычислительная погрешность).
\end{enumerate}

Погрешности делятся на \textbf{неустранимые} (причины 1, 2 --- уменьшаются средствами, внешними по отношению к вычислительной задаче: лучшая модель, более точные измерения) и \textbf{устранимые} (причины 3, 4 --- уменьшаются выбором лучшего метода, увеличением разрядности, рациональной организацией вычислений).

\subsubsection*{Абсолютная и относительная погрешности}

Пусть $a$ --- точное значение, $a^*$ --- приближённое.

\textbf{Погрешность:} $\varepsilon = a - a^*$.

\textbf{Абсолютная погрешность:}
\[
    \Delta a^* = |a - a^*|.
\]

На практике используют \textit{предельную абсолютную погрешность} $\overline{\Delta} a^*$ такую, что $\Delta a^* \leq \overline{\Delta} a^*$:
\[
    a \in \bigl[a^* - \overline{\Delta} a^*,\; a^* + \overline{\Delta} a^*\bigr].
\]

\textbf{Относительная погрешность:}
\[
    \delta a^* = \frac{\Delta a^*}{|a|} \approx \frac{\Delta a^*}{|a^*|}.
\]

Верхняя оценка: $\overline{\delta} a^* \geq \delta a^*$. Обычно выражается в процентах.

\subsubsection*{Значащие цифры в десятичной записи числа}

Приближённое число записывается в виде
\[
    a^* = \pm\, 0.\underbrace{d_1 d_2 \ldots d_m}_{\text{мантисса}} \times 10^p,
\]
где $d_1 \neq 0$.

\textbf{Значащие цифры} --- все цифры мантиссы, начиная с первой ненулевой слева. Нули перед первой ненулевой цифрой --- незначащие. Значащие нули справа убирать нельзя --- они указывают точность (например, $2{,}0$ и $2{,}00$ --- разные приближения).

\subsubsection*{Округление}

\textbf{Округление усечением:} отбрасываются все цифры правее заданного разряда. Абсолютная погрешность не превосходит единицы разряда последней оставленной цифры.

\textbf{Округление по дополнению} (по умолчанию):
\begin{enumerate}
    \item Если первая отбрасываемая цифра $< 5$ --- оставшиеся не меняются.
    \item Если $> 5$ или $= 5$ и среди остальных отбрасываемых есть ненулевые --- последняя оставляемая увеличивается на~1.
    \item Если $= 5$ и все остальные нули --- последняя оставляемая не меняется, если чётная, и увеличивается на~1, если нечётная.
\end{enumerate}

Абсолютная погрешность при округлении по дополнению не превосходит $\frac{1}{2}$ единицы разряда последней оставленной цифры.

\subsubsection*{Верные значащие цифры}

\textbf{Верная в широком смысле:} значащая цифра, для которой абсолютная погрешность не превосходит \textit{единицы} её разряда (связано с округлением усечением).

\textbf{Верная в узком смысле} (по умолчанию): значащая цифра, для которой абсолютная погрешность не превосходит \textit{половины единицы} её разряда (связано с округлением по дополнению).

Верная в узком смысле цифра верна и в широком; обратное неверно.

\textbf{Быстрое правило определения верных цифр:}
\begin{enumerate}
    \item Округлить $\overline{\Delta} a^*$ с избытком до одной значащей цифры.
    \item Если эта цифра $\leq 5$ --- все разряды левее верные.
    \item Если $> 5$ --- соседний слева разряд сомнительный, все левее него --- верные.
\end{enumerate}

\textbf{Обратная задача:} если все цифры верные (в узком смысле), то $\overline{\Delta} a^* = \frac{1}{2}$ единицы разряда последней значащей цифры.

\subsubsection*{Правила записи приближённых чисел}

\begin{enumerate}
    \item В промежуточных результатах оставляют верные $+$ 1--2 сомнительные цифры.
    \item Окончательный результат округляют с сохранением не более одной сомнительной цифры.
\end{enumerate}

\subsubsection*{Пример (практика 1, задача 1)}

Дано число $a = \frac{1}{3}$ и его приближение $a^* = 0{,}3333$. Найти абсолютную и относительную погрешности и определить верные знаки $a^*$.

Абсолютная погрешность:
\[
    \Delta a^* = |a - a^*| = \left|\frac{1}{3} - 0{,}3333\right| = 0{,}0000\overline{3} \approx 0{,}00003.
\]
Относительная погрешность:
\[
    \delta a^* = \frac{\Delta a^*}{|a|} = \frac{0{,}00003}{1/3} \approx 0{,}0001 = 0{,}01\%.
\]
Проверка верных цифр: единицы разрядов всех цифр $a^* = 0{,}3333$ больше абсолютной погрешности ($10^{-1}, 10^{-2}, 10^{-3}, 10^{-4}$ --- все $> 0{,}00003$), поэтому все цифры верные.

\subsubsection*{Пример (лекция 1-1, с.~7, верные цифры по быстрому правилу)}

Пусть $a^* = 3{,}2538$, $\overline{\Delta} a^* = 0{,}042$. Округляем погрешность с избытком до одной значащей цифры: $\overline{\Delta} a^* \approx 0{,}05$. Единственная значащая цифра равна $5$. По правилу ($\leq 5$) все разряды левее верные: верные цифры --- $3$, $2$. Цифра $5$ --- сомнительная.

Пусть $a^* = 0{,}3478$, $\overline{\Delta} a^* = 0{,}082$. Округляем: $\overline{\Delta} a^* \approx 0{,}09$, значащая цифра $9 > 5$. По правилу: соседний слева разряд ($4$) сомнительный, все левее ($3$) --- верные. Итого верная цифра только $3$.


\subsection{Погрешности арифметических операций}

Пусть $a^*$, $b^*$ --- приближённые значения с погрешностями $\Delta a^*$, $\Delta b^*$ и $\delta a^*$, $\delta b^*$.

\textbf{Теорема 1} (абсолютная погрешность суммы и разности):
\[
    \Delta(a^* \pm b^*) \leq \Delta a^* + \Delta b^*.
\]

Обобщение: $\displaystyle\Delta\!\Bigl(\sum_{i=1}^{n} a_i^*\Bigr) \leq \sum_{i=1}^{n} \Delta a_i^*$.

\textbf{Теорема 2} (относительная погрешность суммы и разности, $a, b$ одного знака):
\[
    \delta(a^* + b^*) \leq \max\{\delta a^*,\, \delta b^*\},
\]
\[
    \delta(a^* - b^*) \leq \max\{\delta a^*,\, \delta b^*\} \cdot \frac{|a^*| + |b^*|}{|a^* - b^*|}.
\]

\textit{Следствие:} при вычитании близких чисел ($|a^* - b^*| \to 0$) происходит \textbf{катастрофическая потеря точности}.

\textbf{Теорема 3} (относительная погрешность произведения):
\[
    \delta(a^* \cdot b^*) \leq \delta a^* + \delta b^* + \delta a^* \cdot \delta b^*.
\]

\textit{Следствие:} $\delta(a^* \cdot b^*) \approx \delta a^* + \delta b^*$. Обобщение: $\displaystyle\delta\!\Bigl(\prod_{i=1}^{n} a_i^*\Bigr) \leq \sum_{i=1}^{n} \delta a_i^*$.

\textbf{Теорема 4} (относительная погрешность частного, $b^* \neq 0$, $\delta b^* < 1$):
\[
    \delta\!\left(\frac{a^*}{b^*}\right) \leq \frac{\delta a^* + \delta b^* + \delta a^* \cdot \delta b^*}{1 - \delta b^*}.
\]

\textit{Следствие:} $\displaystyle\delta\!\left(\frac{a^*}{b^*}\right) \approx \delta a^* + \delta b^*$.

\subsubsection*{Пример (практика 1, задача 2)}

Даны $a = \frac{5}{7}$, $b = \frac{3}{11}$ и их приближения $a^* = 0{,}714$, $b^* = 0{,}273$. Тогда:
\[
    \Delta a^* = |a - a^*| = 0{,}000286, \quad \Delta b^* = |b - b^*| = 0{,}000273.
\]
\[
    \delta a^* = \frac{0{,}000286}{5/7} = 0{,}0004, \quad \delta b^* = \frac{0{,}000273}{3/11} = 0{,}001.
\]

Погрешности арифметических операций по определению:
\[
    \Delta(a^* + b^*) = |(a+b) - (a^*+b^*)| = 0{,}000013, \quad
    \Delta(a^* - b^*) = 0{,}000559.
\]

Оценки по теоремам:
\[
    \overline{\Delta}(a^* \pm b^*) = \Delta a^* + \Delta b^* = 0{,}000559.
\]

Относительная погрешность разности по теореме~2:
\[
    \overline{\delta}(a^* - b^*) \leq \max\{0{,}0004;\; 0{,}001\} \cdot \frac{0{,}714 + 0{,}273}{|0{,}714 - 0{,}273|} = 0{,}001 \cdot \frac{0{,}987}{0{,}441} = 0{,}00224.
\]

Относительная погрешность произведения по следствию теоремы~3:
\[
    \overline{\delta}(a^* \cdot b^*) \approx \delta a^* + \delta b^* = 0{,}0004 + 0{,}001 = 0{,}0014.
\]

Относительная погрешность частного по следствию теоремы~4:
\[
    \overline{\delta}\!\left(\frac{a^*}{b^*}\right) \approx \delta a^* + \delta b^* = 0{,}0014.
\]

\subsubsection*{Пример (практика 1, задача 3)}

Даны $a^* = 4{,}362$, $b^* = 2{,}743$, все знаки верные. Вычислить
\[
    C^* = \frac{a^* \cdot (b^* - 2)}{a^* \cdot b^*}.
\]

Все цифры верные $\Rightarrow$ $\Delta a^* = 0{,}0005$, $\Delta b^* = 0{,}0005$.

$\delta a^* = 0{,}0005/4{,}362 = 0{,}000115$, $\delta b^* = 0{,}0005/2{,}743 = 0{,}000182$.

В разности $b^* - 2$ константа $2$ задана точно, поэтому $\delta(b^*-2) = \delta b^* \cdot \frac{|b^*|}{|b^*-2|} = 0{,}000182 \cdot \frac{2{,}743}{0{,}743} = 0{,}000672$.

Относительная погрешность числителя (произведение): $\delta \approx \delta a^* + \delta(b^*-2) = 0{,}000787$.

Относительная погрешность знаменателя (произведение): $\delta \approx \delta a^* + \delta b^* = 0{,}000297$.

Обозначим $A = a^*(b^*-2) = 3{,}241$, $B = a^* \cdot b^* = 11{,}965$. Тогда
\[
    C^* = A - B = 3{,}241 - 11{,}965 = -8{,}724.
\]

Оба числа одного знака при вычитании (оба положительных, берём разность по модулю), относительная погрешность разности:
\[
    \delta C^* \leq \max\{0{,}000787;\; 0{,}000297\} \cdot \frac{3{,}241 + 11{,}965}{|3{,}241 - 11{,}965|} = 0{,}000787 \cdot \frac{15{,}206}{8{,}724} = 0{,}00137.
\]

$\overline{\Delta} C^* = 0{,}00137 \cdot 8{,}724 = 0{,}012$. Верная цифра --- до разряда десятых; с одной сомнительной: $C^* \approx -8{,}72$.


\subsection{Погрешности функций. Линейная оценка дифференциальной погрешности}

Пусть $y = f(x_1, \ldots, x_n)$ --- непрерывно дифференцируемая функция, $x_i^*$ --- приближённые значения аргументов с предельными погрешностями $\Delta x_i^*$, $\delta x_i^*$.

\textbf{Линейная оценка абсолютной погрешности:}
\[
    \Delta y^* \approx \sum_{i=1}^{n} \left|\frac{\partial f}{\partial x_i}(x_1^*, \ldots, x_n^*)\right| \cdot \Delta x_i^*.
\]

\textbf{Оценка относительной погрешности:}
\[
    \delta y^* \approx \sum_{i=1}^{n} \left|\frac{x_i^* \cdot f'_{x_i}(x_1^*, \ldots, x_n^*)}{f(x_1^*, \ldots, x_n^*)}\right| \cdot \delta x_i^*.
\]

\textbf{Частные случаи:}
\begin{itemize}
    \item $f = x_1 \pm x_2$: $\;\Delta f^* = \Delta x_1^* + \Delta x_2^*$ (теорема~1).
    \item $f = x_1^{p_1} \cdot x_2^{p_2} \cdots x_n^{p_n}$: $\;\delta f^* = \sum |p_i|\,\delta x_i^*$ (обобщение теорем~3,~4).
\end{itemize}

\textbf{Обратная задача} (принцип равного влияния): если требуется $\Delta y^* \leq \varepsilon$, то вклад каждого аргумента полагаем равным:
\[
    \left|\frac{\partial f}{\partial x_i}\right| \cdot \Delta x_i^* = \frac{\varepsilon}{n}
    \quad\Longrightarrow\quad
    \Delta x_i^* = \frac{\varepsilon}{\displaystyle n\,\left|\frac{\partial f}{\partial x_i}\right|}.
\]

\subsubsection*{Пример (практика 1, задача 4)}

Даны $a^* = 4{,}362$, $b^* = 2{,}743$, все знаки верные. Рассмотреть функцию
\[
    f(a, b) = a^2 \cdot b - a \cdot b^2.
\]

$\Delta a^* = 0{,}0005$, $\Delta b^* = 0{,}0005$. Вычисляем значение: $f^* = 4{,}362^2 \cdot 2{,}743 - 4{,}362 \cdot 2{,}743^2$.

Частные производные:
\[
    \frac{\partial f}{\partial a} = 2ab - b^2, \quad
    \frac{\partial f}{\partial b} = a^2 - 2ab.
\]

При $a^* = 4{,}362$, $b^* = 2{,}743$:
\[
    \left|\frac{\partial f}{\partial a}\right| = |2 \cdot 4{,}362 \cdot 2{,}743 - 2{,}743^2| = |23{,}922 - 7{,}524| = 16{,}398,
\]
\[
    \left|\frac{\partial f}{\partial b}\right| = |4{,}362^2 - 2 \cdot 4{,}362 \cdot 2{,}743| = |19{,}027 - 23{,}922| = 4{,}895.
\]

Абсолютная погрешность:
\[
    \overline{\Delta} f^* = 16{,}398 \cdot 0{,}0005 + 4{,}895 \cdot 0{,}0005 = 0{,}0082 + 0{,}0024 = 0{,}011.
\]

$f^* = 52{,}178 - 32{,}821 = 19{,}357$. Относительная погрешность: $\overline{\delta} f^* = 0{,}011/19{,}357 = 0{,}0006$. С одной сомнительной цифрой: $f^* \approx 19{,}36$.

\subsubsection*{Пример (практика 1, задача 5, обратная задача)}

Функция $f(a,b) = a^2 \cdot b - a \cdot b^2$, $a^* = 4{,}362$, $b^* = 2{,}743$. С какой точностью надо задать $a^*$, $b^*$, чтобы $\overline{\Delta} f^* \leq 0{,}01$?

По принципу равных влияний ($n = 2$):
\[
    \Delta a^* = \frac{0{,}01}{2 \cdot |\partial f/\partial a|} = \frac{0{,}01}{2 \cdot 16{,}398} = 0{,}0003,
\]
\[
    \Delta b^* = \frac{0{,}01}{2 \cdot |\partial f/\partial b|} = \frac{0{,}01}{2 \cdot 4{,}895} = 0{,}001.
\]

Значит, $a^*$ нужно знать с точностью $0{,}0003$ (4 верных знака: $a^* = 4{,}362$), а $b^*$ --- с точностью $0{,}001$ (3 верных знака: $b^* = 2{,}743$).

\subsubsection*{Пример (лекция 1-1, с.~20--21)}

Найти значение функции $y = \sqrt{\frac{a}{b}}$, если $a^* = 1{,}23$, $b^* = 0{,}456$. Оценить погрешности.

$\Delta a^* = 0{,}005$, $\Delta b^* = 0{,}0005$. $\delta a^* = 0{,}005/1{,}23 = 0{,}004$, $\delta b^* = 0{,}0005/0{,}456 = 0{,}001$.

$y^* = \sqrt{1{,}23/0{,}456} = \sqrt{2{,}697} = 1{,}642$.

\textbf{Способ 1} (линейная оценка). Частные производные:
\[
    \frac{\partial y}{\partial a} = \frac{1}{2}\sqrt{\frac{1}{ab}} = \frac{1}{2\sqrt{ab}}, \quad
    \frac{\partial y}{\partial b} = -\frac{1}{2}\sqrt{\frac{a}{b^3}} = -\frac{\sqrt{a}}{2b\sqrt{b}}.
\]
\[
    \left|\frac{\partial y}{\partial a}\right| = \frac{1}{2\sqrt{1{,}23 \cdot 0{,}456}} = 0{,}667, \quad
    \left|\frac{\partial y}{\partial b}\right| = \frac{\sqrt{1{,}23}}{2 \cdot 0{,}456 \cdot \sqrt{0{,}456}} = 1{,}800.
\]
\[
    \Delta y^* = 0{,}667 \cdot 0{,}005 + 1{,}800 \cdot 0{,}0005 = 0{,}003 + 0{,}001 = 0{,}004.
\]
$\delta y^* = 0{,}004 / 1{,}642 = 0{,}002$.

\textbf{Способ 2} (через относительную погрешность). Для $y = a^{1/2} \cdot b^{-1/2}$:
\[
    \delta y^* = \tfrac{1}{2}\,\delta a^* + \tfrac{1}{2}\,\delta b^* = \tfrac{1}{2} \cdot 0{,}004 + \tfrac{1}{2} \cdot 0{,}001 = 0{,}0025.
\]
$\Delta y^* = 0{,}0025 \cdot 1{,}642 = 0{,}004$.

Оба способа: $y^* \approx 1{,}64$ (с одной сомнительной цифрой), $\overline{\Delta} y^* = 0{,}004$, $\overline{\delta} y^* = 0{,}002$.


\newpage
\section{Часть II. Вычислительные задачи и методы}

\subsection{Особенности современных вычислительных задач. Этапы решения вычислительной задачи. Вычислительные задачи, их корректность и обусловленность. Классификация вычислительных методов}

\subsubsection*{Особенности современных вычислительных задач}

Практически все числовые величины при расчётах являются приближёнными. Вычислительная задача --- это задание входных числовых данных и формулировка требуемого числового результата. Современные задачи характеризуются большой размерностью, сложностью моделей и необходимостью контроля погрешностей на каждом этапе.

\subsubsection*{Этапы решения вычислительной задачи}

\begin{enumerate}
    \item Построение математической модели.
    \item Выбор метода решения.
    \item Разработка алгоритма.
    \item Программирование и отладка.
    \item Проведение расчётов.
    \item Анализ результатов.
\end{enumerate}

\subsubsection*{Корректность вычислительной задачи}

Решение \textit{устойчиво}, если малым возмущениям исходных данных соответствуют малые погрешности решения.

Формально (скалярный случай): $\forall\,\varepsilon > 0\;\;\exists\,\delta(\varepsilon) > 0$: если $\Delta x^* < \delta(\varepsilon)$, то $\Delta y^* < \varepsilon$.

\textbf{Вычислительная задача корректна}, если:
\begin{enumerate}
    \item Для любых допустимых исходных данных существует единственное решение.
    \item Решение устойчиво по отношению к малым возмущениям исходных данных.
\end{enumerate}

При нарушении хотя бы одного условия задача \textit{некорректна}.

\subsubsection*{Обусловленность}

\textbf{Число обусловленности} --- коэффициент возможного роста погрешности решения по отношению к погрешностям входных данных:
\[
    \Delta y^* \leq \vartheta_\Delta \cdot \Delta x^*, \qquad
    \delta y^* \leq \vartheta_\delta \cdot \delta x^*.
\]

Для задачи вычисления функции $y = f(x)$:
\[
    \vartheta_\Delta = |f'(x^*)|, \qquad
    \vartheta_\delta = \left|\frac{x^* \cdot f'(x^*)}{y^*}\right|.
\]

Если $\vartheta$ мало --- задача \textit{хорошо обусловлена}; если $\vartheta \gg 1$ --- \textit{плохо обусловлена}.

\subsubsection*{Пример (лекция 1-2, плохо обусловленная система)}

Рассмотрим систему
\[
    \begin{cases}
        1{,}1\, x_1 + x_2 = 1{,}1, \\
        x_1 + x_2 = 1,
    \end{cases}
\]
имеющую решение $x_1 = 1$, $x_2 = 0$. Изменим коэффициент во втором уравнении: $(1+\varepsilon)\,x_1 + x_2 = 1$. Тогда решение:
\[
    x_1 = \frac{1}{1 - 10\varepsilon}, \quad x_2 = \frac{-11\varepsilon}{1 - 10\varepsilon}.
\]

При $\varepsilon = 0{,}001$: $x_1 = 1{,}01$, $x_2 = -0{,}011$. При $\varepsilon = 0{,}002$: $x_1 = 1{,}02$, $x_2 = -0{,}022$.

Считая $\varepsilon$ входным данным: $\Delta\varepsilon^* = 0{,}001$, $\Delta x_2^* = 0{,}011$, поэтому $\vartheta_\Delta = 0{,}011/0{,}001 = 11$. Абсолютное число обусловленности велико --- задача плохо обусловлена.

\subsubsection*{Классификация вычислительных методов}

\begin{enumerate}
    \item \textbf{Методы эквивалентных преобразований.} Исходная задача заменяется эквивалентной с тем же решением, но более простой. Пример: $f(x) = 0 \;\Leftrightarrow\; \min g(x) = (f(x))^2$.

    \item \textbf{Методы аппроксимации.} Исходная задача заменяется другой, решение которой близко к решению исходной. Важна оценка погрешности аппроксимации. Пример: вычисление интеграла через интерполяционный многочлен.

    \item \textbf{Прямые методы.} Точное решение за конечное число шагов. Примеры: формулы корней, метод Гаусса, метод прогонки.

    \item \textbf{Итерационные методы.} Строится последовательность приближений $\{y_k\}$, сходящаяся к точному решению. Метод \textit{одношаговый}, если каждая итерация использует одну предыдущую, и \textit{многошаговый}, если несколько. Метод \textit{сходится}, если последовательность сходится к точному решению; множество начальных приближений, обеспечивающих сходимость --- \textit{область сходимости}.

    Для итерационного метода необходимы:
    \begin{itemize}
        \item описание итерационного шага (расчётная формула);
        \item область/условия сходимости;
        \item оценка погрешности для критерия останова.
    \end{itemize}

    Примеры: метод простой итерации, метод Ньютона, методы Якоби и Зейделя.
\end{enumerate}

\subsubsection*{Пример итерационного метода (лекция 1-2, метод Герона)}

Задача: вычислить $\sqrt{x}$, т.е. найти $a$ такое, что $a^2 = x$. Итерационная формула:
\[
    a_{k+1} = \frac{1}{2}\left(a_k + \frac{x}{a_k}\right).
\]

Для $x = 2$ с начальным приближением $a_0 = 1$:
\[
    a_1 = \frac{1}{2}\left(1 + 2\right) = 1{,}5, \quad
    a_2 = \frac{1}{2}\left(1{,}5 + \frac{2}{1{,}5}\right) = 1{,}41667, \quad
    a_3 = 1{,}41421.
\]

Все значащие цифры $a_3$ верные ($\sqrt{2} = 1{,}41421\ldots$). Сходится при любом положительном начальном приближении.

\subsubsection*{Корректность алгоритмов}

\textbf{Вычислительный алгоритм корректен}, если:
\begin{enumerate}
    \item За конечное число операций преобразует любое допустимое $x \in X$ в единственный результат $y \in Y$.
    \item Результат устойчив по входным данным (при отсутствии вычислительной погрешности).
    \item Результат вычислительно устойчив (вычислительная погрешность $\to 0$ при машинной погрешности $\to 0$).
\end{enumerate}

Алгоритм \textit{хорошо обусловлен}, если $\delta y^* \leq \vartheta_A \cdot \varepsilon_M$ при малом $\vartheta_A$, где $\varepsilon_M$ --- относительная погрешность округления, $\vartheta_A$ --- число обусловленности алгоритма.

\section{Часть III. Аналитическое приближение табличных функций}

\subsection{Задача приближенного вычисления функции. Интерполяция. Задача полиномиальной интерполяции. Теорема о единственности решения. Погрешность полиномиальной интерполяции}

\subsubsection*{Задача приближенного вычисления функции}

Пусть функция $y = f(x)$ задана таблицей значений $\{(x_i, y_i)\}_{i=0}^{n}$ или вычисляется по сложной формуле. Задача: построить просто вычисляемую функцию $g(x)$ (например, полином), приближённо равную $f$. Два подхода: \textit{интерполяция} (значения $g$ совпадают с табличными в узлах) и \textit{аппроксимация} (минимизируется некоторая мера отклонения).

\subsubsection*{Интерполяция}

По таблице $\{(x_i, y_i)\}_{i=0}^{n}$ строится функция $g(x)$, удовлетворяющая условиям интерполяции:
\[
    g(x_i) = y_i, \quad i = 0, 1, \ldots, n.
\]

Вне узлов полагаем $f(x) \approx g(x)$. Функция $g$ --- интерполирующая, точки $x_0, \ldots, x_n$ --- узлы интерполяции. Для единственности решения задаётся класс интерполирующих функций $\mathfrak{G}$.

\subsubsection*{Задача полиномиальной интерполяции}

Класс $\mathfrak{G}$ --- алгебраические полиномы степени $m$:
\[
    P_m(x) = a_0 + a_1 x + \cdots + a_m x^m.
\]

Особый интерес представляет случай $m = n$ (степень многочлена на единицу меньше числа узлов).

\subsubsection*{Теорема о единственности решения}

\textbf{Теорема 1.} Если узлы $x_0, x_1, \ldots, x_n$ попарно различны, то существует единственный интерполяционный многочлен $P_n$ степени $n$.

\textit{Идея:} условия интерполяции $P_n(x_i) = y_i$ дают систему $n+1$ линейного уравнения с $n+1$ неизвестным. Определитель системы --- определитель Вандермонда, который отличен от нуля при попарно различных узлах.

При $m > n$ --- бесконечно много решений; при $m < n$ --- решений может не быть.

\subsubsection*{Погрешность полиномиальной интерполяции}

\textbf{Теорема 2} (об остаточном члене). Пусть $f$ $(n+1)$ раз непрерывно дифференцируема на $[a; b] \supset \{x_0, \ldots, x_n\}$. Тогда $\forall\, x \in [a; b]$ $\exists\, \zeta \in (a; b)$:
\[
    R_n(x) = f(x) - P_n(x) = \frac{f^{(n+1)}(\zeta)}{(n+1)!}\, Q_{n+1}(x),
\]
где $Q_{n+1}(x) = (x - x_0)(x - x_1) \cdots (x - x_n)$.

Оценочная функция погрешности:
\[
    \overline{\Delta}\bigl(P_n(x)\bigr) = \frac{M_{n+1}}{(n+1)!}\, |Q_{n+1}(x)|, \quad M_{n+1} = \max_{z \in [a;b]} |f^{(n+1)}(z)|.
\]

Глобальная оценка: $V(P_n) = \sup_{x \in [a;b]} \overline{\Delta}(P_n(x))$.

\subsubsection*{Пример (практика 2, задача 1)}

Дана функция $f(x) = \frac{1}{\sqrt{x}}$, начальный узел $x_0 = 0{,}5$, шаг $h = 0{,}25$, три узла ($n = 2$). Оценить погрешность интерполяции.

Узлы: $x_0 = 0{,}5$, $x_1 = 0{,}75$, $x_2 = 1{,}0$. Можно построить полином степени 2. Третья производная:
\[
    f'''(x) = -\frac{15}{8}\, x^{-7/2}.
\]
На отрезке $[0{,}5;\; 1{,}0]$ максимум $|f'''|$ достигается при $x = 0{,}5$: $M_3 = |f'''(0{,}5)| \approx 33{,}94$.

Оценка:
\[
    \overline{\Delta}(P_2(x)) = \frac{M_3}{3!}\,|Q_3(x)| = \frac{33{,}94}{6}\,|(x - 0{,}5)(x - 0{,}75)(x - 1{,}0)|.
\]
Исследованием $Q_3(x)$ на $[0{,}5;\; 1]$ находим $\max |Q_3| \approx 0{,}0048$. Наибольшая погрешность: $V(P_2) \approx 33{,}94 \cdot 0{,}0048 / 6 \approx 0{,}027$.

\subsubsection*{Пример (практика 2, задача 2)}

С какой точностью можно вычислить $\sqrt{5}$ по известным $\sqrt{1}$, $\sqrt{4}$, $\sqrt{9}$, $\sqrt{16}$?

Число узлов $n + 1 = 4$, $n = 3$. Функция $f(x) = \sqrt{x}$, четвёртая производная:
\[
    f^{(4)}(x) = \frac{15}{16}\, x^{-7/2}.
\]
$|f^{(4)}(x)|$ убывает; максимум на $[1; 16]$: $M_4 = f^{(4)}(1) = 15/16$.

$Q_4(5) = (5-1)(5-4)(5-9)(5-16) = 4 \cdot 1 \cdot (-4) \cdot (-11) = 176$.

Погрешность:
\[
    |R_3(5)| \leq \frac{M_4}{4!}\, |Q_4(5)| = \frac{15/16}{24} \cdot 176 \approx 0{,}69.
\]
Итого $\sqrt{5}$ можно вычислить с одним верным знаком после запятой.


\subsection{Интерполяционный многочлен Лагранжа. Схема Эйткена}

\subsubsection*{Интерполяционный многочлен Лагранжа}

Для таблицы $\{(x_i, y_i)\}_{i=0}^{n}$ с попарно различными узлами:
\[
    L_n(x) = \sum_{i=0}^{n} y_i \cdot l_{n,i}(x), \quad l_{n,i}(x) = \prod_{\substack{j=0\\j\neq i}}^{n} \frac{x - x_j}{x_i - x_j}.
\]

Величины $l_{n,i}(x)$ --- множители Лагранжа.

Частные случаи: при $n=1$ --- прямая через две точки; при $n=2$ --- парабола через три точки.

Недостаток: при добавлении нового узла формулу надо пересчитывать заново.

\subsubsection*{Схема Эйткена}

Рекуррентное вычисление значения $L_n(x)$ в точке $x$. Обозначим $L_k^{i,\ldots,i+k}(x)$ --- значение многочлена Лагранжа степени $k$ по узлам $x_i, \ldots, x_{i+k}$.

Начальные значения: $L_0^i(x) = y_i$.

Рекуррентная формула $k$-го шага:
\[
    L_k^{i,\ldots,i+k}(x) = \frac{1}{x_{i+k} - x_i}
    \begin{vmatrix}
        L_{k-1}^{i,\ldots,i+k-1}(x) & x_i - x \\
        L_{k-1}^{i+1,\ldots,i+k}(x) & x_{i+k} - x
    \end{vmatrix}.
\]

На последнем шаге получается $L_n(x) = L_n^{0,\ldots,n}(x)$.

\subsubsection*{Пример (практика 2, задача 3)}

Таблица:

\begin{center}
\begin{tabular}{c|cccc}
    $x_i$ & $-1$ & $0$ & $1$ & $3$ \\
    \hline
    $y_i$ & $12$ & $12$ & $0$ & $60$
\end{tabular}
\end{center}

Многочлен Лагранжа:
\[
    L_3(x) = 12 \cdot \frac{x(x-1)(x-3)}{(-1)(-2)(-4)} + 12 \cdot \frac{(x+1)(x-1)(x-3)}{(1)(-1)(-3)} + 0 + 60 \cdot \frac{(x+1)x(x-1)}{(4)(3)(2)}.
\]
После упрощения: $L_3(x) = 5x^3 - 6x^2 - 11x + 12$.

\subsubsection*{Пример (практика 2, задача 5, схема Эйткена)}

Таблица:

\begin{center}
\begin{tabular}{c|cccccc}
    $x_i$ & $-2{,}0$ & $-1{,}5$ & $-1{,}0$ & $-0{,}5$ & $0{,}0$ & $0{,}5$ \\
    \hline
    $f(x_i)$ & $3{,}2$ & $3{,}7$ & $3{,}3$ & $2{,}8$ & $2{,}5$ & $2{,}0$
\end{tabular}
\end{center}

Вычислить $f(-1{,}3)$ по схеме Эйткена. Третий столбец таблицы --- разности $x_i - x$, где $x = -1{,}3$. На первом шаге:
\[
    L_1^{0,1}(-1{,}3) = \frac{1}{-1{,}5+2}
    \begin{vmatrix}
    3{,}2 & -0{,}7 \\
    3{,}7 & -0{,}2
    \end{vmatrix}
    = 2(-0{,}2 \cdot 3{,}2 + 0{,}7 \cdot 3{,}7) = 3{,}9.
\]

Продолжая рекуррентно, получаем $f(-1{,}3) \approx L_5(-1{,}3) = 3{,}59$.


\subsection{Разделенные разности и их свойства. Интерполяционные многочлены Ньютона с разделенными разностями}

\subsubsection*{Разделённые разности}

\textbf{Первого порядка:}
\[
    f(x_i;\, x_{i+1}) = \frac{f(x_{i+1}) - f(x_i)}{x_{i+1} - x_i}.
\]

\textbf{$k$-го порядка} (рекуррентно):
\[
    f(x_i;\, \ldots;\, x_{i+k}) = \frac{f(x_{i+1};\, \ldots;\, x_{i+k}) - f(x_i;\, \ldots;\, x_{i+k-1})}{x_{i+k} - x_i}.
\]

\textbf{Свойства:}
\begin{itemize}
    \item Разделённые разности --- симметричные функции своих аргументов (не зависят от порядка узлов).
    \item Связь с конечными разностями (для равномерной таблицы с шагом $h$):
    $f(x_i;\, \ldots;\, x_{i+k}) = \dfrac{\Delta^k y_i}{k!\, h^k}$.
\end{itemize}

\subsubsection*{Многочлен Ньютона I вида (вперёд)}

Строится по верхней диагонали таблицы разделённых разностей:
\[
    P_n^{(1)}(x) = f(x_0) + \sum_{i=1}^{n} f(x_0;\, \ldots;\, x_i)\,(x - x_0)(x - x_1) \cdots (x - x_{i-1}).
\]

Применяется для интерполяции в начале таблицы.

\subsubsection*{Многочлен Ньютона II вида (назад)}

Строится по нижней диагонали:
\[
    P_n^{(2)}(x) = f(x_n) + \sum_{i=1}^{n} f(x_{n-i};\, \ldots;\, x_n)\,(x - x_n)(x - x_{n-1}) \cdots (x - x_{n-i+1}).
\]

Применяется для интерполяции в конце таблицы.

Преимущество перед формулой Лагранжа: при добавлении нового узла достаточно вычислить и прибавить одно слагаемое.

Приближённая оценка погрешности:
\[
    \overline{\Delta}(P_n(x)) \approx |f(x;\, x_0;\, \ldots;\, x_n) \cdot Q_{n+1}(x)|.
\]

\subsubsection*{Пример (практика 2, задача 6)}

Та же таблица: $(-1, 12)$, $(0, 12)$, $(1, 0)$, $(3, 60)$.

Таблица разделённых разностей:

\begin{center}
\begin{tabular}{c|c|c|c|c}
    $x_i$ & $f(x_i)$ & $f(x_i; x_{i+1})$ & $f(\ldots)$ II пор. & $f(\ldots)$ III пор. \\
    \hline
    $-1$ & $12$ & & & \\
         &     & $0$ & & \\
    $0$  & $12$ &    & $-6$ & \\
         &     & $-12$ &   & $5$ \\
    $1$  & $0$  &    & $14$ & \\
         &     & $30$ & & \\
    $3$  & $60$ & & &
\end{tabular}
\end{center}

Многочлен I вида (верхняя диагональ: $12,\; 0,\; -6,\; 5$):
\[
    P_3^{(1)}(x) = 12 + 0 \cdot (x+1) - 6(x+1)x + 5(x+1)x(x-1) = 5x^3 - 6x^2 - 11x + 12.
\]

Многочлен II вида (нижняя диагональ: $60,\; 30,\; 14,\; 5$):
\[
    P_3^{(2)}(x) = 60 + 30(x-3) + 14(x-3)(x-1) + 5(x-3)(x-1)x = 5x^3 - 6x^2 - 11x + 12.
\]

Результаты совпадают --- единственный интерполяционный многочлен степени 3.


\subsection{Конечные разности и их свойства. Интерполяционные многочлены Ньютона с конечными разностями}

\subsubsection*{Конечные разности}

Для равномерной таблицы ($x_i = x_0 + ih$, $i = 0, \ldots, n$).

\textbf{Первого порядка:} $\Delta y_i = y_{i+1} - y_i$.

\textbf{$k$-го порядка:} $\Delta^k y_i = \Delta^{k-1} y_{i+1} - \Delta^{k-1} y_i$.

Явная формула через биномиальные коэффициенты:
\[
    \Delta^k y_i = \sum_{j=0}^{k} (-1)^j \binom{k}{j}\, y_{i+k-j}.
\]

Связь с разделёнными разностями: $f(x_i;\, \ldots;\, x_{i+k}) = \dfrac{\Delta^k y_i}{k!\, h^k}$.

\subsubsection*{Многочлен Ньютона I вида с конечными разностями}

Замена $q = \dfrac{x - x_0}{h}$:
\[
    P_n^{(1)}(x_0 + qh) = y_0 + \sum_{i=1}^{n} \frac{\Delta^i y_0}{i!}\, q(q-1) \cdots (q-i+1).
\]

Разности берутся из верхней строки таблицы конечных разностей.

\subsubsection*{Многочлен Ньютона II вида с конечными разностями}

Замена $q = \dfrac{x - x_n}{h}$:
\[
    P_n^{(2)}(x_n + qh) = y_n + \sum_{i=1}^{n} \frac{\Delta^i y_{n-i}}{i!}\, q(q+1) \cdots (q+i-1).
\]

Разности --- из нижней строки таблицы.

Первый многочлен лучше для начала таблицы, второй --- для конца.

\subsubsection*{Пример (практика 2, задача 8)}

Таблица:

\begin{center}
\begin{tabular}{c|cccc}
    $x_i$ & $0$ & $1$ & $2$ & $3$ \\
    \hline
    $y_i$ & $1$ & $3$ & $7$ & $15$
\end{tabular}
\end{center}

Шаг $h = 1$. Таблица конечных разностей:

\begin{center}
\begin{tabular}{c|c|c|c|c}
    $x_i$ & $y_i$ & $\Delta y_i$ & $\Delta^2 y_i$ & $\Delta^3 y_i$ \\
    \hline
    $0$ & $1$ & & & \\
        &     & $2$ & & \\
    $1$ & $3$ &    & $2$ & \\
        &     & $4$ &   & $2$ \\
    $2$ & $7$ &    & $4$ & \\
        &     & $8$ & & \\
    $3$ & $15$ & & &
\end{tabular}
\end{center}

Многочлен I вида ($q = \frac{x - 0}{1} = x$, верхняя строка: $\Delta y_0 = 2$, $\Delta^2 y_0 = 2$, $\Delta^3 y_0 = 2$):
\[
    P_3^{(1)}(x) = 1 + \frac{2}{1!}\,x + \frac{2}{2!}\,x(x-1) + \frac{2}{3!}\,x(x-1)(x-2) = \frac{x^3 + 2x + 3}{3}.
\]

Многочлен II вида ($q = \frac{x - 3}{1}$, нижняя строка: $\Delta y_2 = 8$, $\Delta^2 y_1 = 4$, $\Delta^3 y_0 = 2$):
\[
    P_3^{(2)}(x) = 15 + \frac{8}{1!}\,q + \frac{4}{2!}\,q(q+1) + \frac{2}{3!}\,q(q+1)(q+2),
\]
что после подстановки $q = x - 3$ даёт тот же результат.


\subsection{Интерполяционные многочлены Гаусса, Стирлинга и Бесселя}

Применяются для интерполяции \textit{в середине} равномерной таблицы с шагом $h$. Центральный узел $x_0 = a$.

\subsubsection*{Первая формула Гаусса}

Узлы подключаются в порядке: $a$, $a+h$, $a-h$, $a+2h$, $a-2h$, \ldots

Замена $q = \dfrac{x - a}{h}$:
\[
    P_n(a + qh) = y_0 + \frac{\Delta y_0}{1!}\,q + \frac{\Delta^2 y_{-1}}{2!}\,q(q-1) + \frac{\Delta^3 y_{-1}}{3!}\,q(q-1)(q+1) + \cdots
\]

Разности идут зигзагом от $y_0$ вниз-вправо в таблице конечных разностей. Применяется при $q > 0$ (точка правее $a$).

\subsubsection*{Вторая формула Гаусса}

Узлы: $a$, $a-h$, $a+h$, $a-2h$, \ldots Замена та же.
\[
    P_n(a + qh) = y_0 + \frac{\Delta y_{-1}}{1!}\,q + \frac{\Delta^2 y_{-1}}{2!}\,q(q+1) + \frac{\Delta^3 y_{-2}}{3!}\,q(q+1)(q-1) + \cdots
\]

Применяется при $q < 0$ (точка левее $a$).

\subsubsection*{Формула Стирлинга}

Получается из первой формулы Гаусса группировкой попарных слагаемых и симметризацией. Число узлов нечётно ($n$ чётно). Применяется при $|q| < \frac{1}{4}$ (точка близко к $a$):
\[
    P_n(a + qh) = y_0 + \frac{q}{1!} \cdot \frac{\Delta y_0 + \Delta y_{-1}}{2} + \frac{q^2}{2!}\,\Delta^2 y_{-1} + \frac{q(q^2 - 1)}{3!} \cdot \frac{\Delta^3 y_{-1} + \Delta^3 y_{-2}}{2} + \cdots
\]

Содержит полусуммы конечных разностей нечётных порядков и разности чётных порядков.

\subsubsection*{Формула Бесселя}

Число узлов чётно ($n$ нечётно). Применяется, когда точка $x$ близка к середине интервала $(a;\, a+h)$, т.е. $q \approx \frac{1}{2}$:
\[
    P_n(a + qh) = \frac{y_0 + y_1}{2} + \frac{q - \frac{1}{2}}{1!}\,\Delta y_0 + \frac{q(q-1)}{2!} \cdot \frac{\Delta^2 y_{-1} + \Delta^2 y_0}{2} + \cdots
\]

Содержит полусуммы разностей чётных порядков и разности нечётных порядков.

При $q = \frac{1}{2}$ формула упрощается --- все слагаемые с нечётными разностями обнуляются.

\subsubsection*{Пример (практика 2, задача 10, формула Бесселя)}

Функция $f(x) = x^3$ на $[0;\, 2]$ с шагом $h = 0{,}4$, точка $x = 1{,}5$, полином степени 4.

Центральный узел $a = x_3 = 1{,}2$, тогда $q = \frac{1{,}5 - 1{,}2}{0{,}4} = 0{,}75$.

Число узлов 5 (чётно), $n = 4$ (нечётно) --- подходит формула Бесселя. Но поскольку $q = 0{,}75$ близко к 1, а не к $\frac{1}{2}$, лучше сдвинуть: $a = x_3 = 1{,}2$ и при $q = 0{,}5$ точка попадёт в $x = 1{,}4$. Для $x = 1{,}5$ берём $a = 1{,}2$, $q = 0{,}75$.

Если же $q = 0{,}5$ (точка --- ровно середина между узлами), формула Бесселя упрощается: остаются только слагаемые с полусуммами чётных разностей. В задаче 10 из практики: $f(1{,}5) = P_4(1{,}5) = 3{,}375$.

\subsection{Задача аппроксимации табличной функции. Аппроксимация в классе обобщённых многочленов методом наименьших квадратов. Полиномиальная аппроксимация}

\subsubsection*{Задача аппроксимации табличной функции}

Пусть функция $y = f(x)$ задана таблицей $\{(x_i, y_i)\}_{i=1}^{n}$. Требуется построить функцию $g(x) \in \mathfrak{G}$, аппроксимирующую исходную, т.е. удовлетворяющую условиям
\[
    g(x_i) \approx y_i, \quad i = 1, \ldots, n.
\]
Точки $x_i$ называются узлами аппроксимации, $\mathfrak{G}$ --- класс аппроксимирующих функций.

В отличие от интерполяции, аппроксимирующая функция не обязана проходить точно через узлы --- она сглаживает данные. Применение интерполяции не всегда оправдано, особенно при большом числе узлов: интерполяционные многочлены имеют высокие степени и вычисляются с большими погрешностями; при наличии ошибок в данных (измерения, статистика) погрешности переходят в интерполирующую функцию. Аппроксимация сглаживает эти погрешности.

\subsubsection*{Метод наименьших квадратов (МНК)}

Аппроксимирующая функция $g \in \mathfrak{G}$ строится из условия минимума среднеквадратического отклонения $g$ от $f$ в узлах:
\[
    \delta_2(g) = \sqrt{\frac{1}{n} \sum_{i=1}^{n} \bigl(y_i - g(x_i)\bigr)^2} \longrightarrow \min.
\]

Задача МНК: по таблице построить функцию $g \in \mathfrak{G}$, для которой $\delta_2(g)$ минимальна.

\subsubsection*{Аппроксимация в классе обобщённых многочленов}

Пусть $\mathfrak{G}$ --- множество обобщённых многочленов степени $m$ по системе линейно независимых базисных функций $\{\varphi_0, \varphi_1, \ldots, \varphi_m\}$:
\[
    \Phi_m(x) = \sum_{j=0}^{m} a_j \varphi_j(x).
\]

Подставляя в $\delta_2$, получаем эквивалентную задачу минимизации:
\[
    S(a_0, \ldots, a_m) = \sum_{i=1}^{n} \left(y_i - \sum_{j=0}^{m} a_j \varphi_j(x_i)\right)^2 \longrightarrow \min_{a_0, \ldots, a_m}.
\]

Приравнивая частные производные к нулю:
\[
    \frac{\partial S}{\partial a_k} = -2 \sum_{i=1}^{n} \left(y_i - \sum_{j=0}^{m} a_j \varphi_j(x_i)\right) \varphi_k(x_i) = 0, \quad k = 0, \ldots, m,
\]
получаем \textbf{нормальную систему МНК}:
\[
    \sum_{j=0}^{m} a_j \sum_{i=1}^{n} \varphi_j(x_i)\,\varphi_k(x_i) = \sum_{i=1}^{n} y_i\,\varphi_k(x_i), \quad k = 0, \ldots, m.
\]

В развёрнутом виде:
\[
    \begin{cases}
        \displaystyle a_0 \sum_{i=1}^{n} \bigl(\varphi_0(x_i)\bigr)^2 + a_1 \sum_{i=1}^{n} \varphi_1(x_i)\,\varphi_0(x_i) + \cdots + a_m \sum_{i=1}^{n} \varphi_m(x_i)\,\varphi_0(x_i) = \sum_{i=1}^{n} y_i\,\varphi_0(x_i), \\[6pt]
        \displaystyle a_0 \sum_{i=1}^{n} \varphi_0(x_i)\,\varphi_1(x_i) + a_1 \sum_{i=1}^{n} \bigl(\varphi_1(x_i)\bigr)^2 + \cdots + a_m \sum_{i=1}^{n} \varphi_m(x_i)\,\varphi_1(x_i) = \sum_{i=1}^{n} y_i\,\varphi_1(x_i), \\[6pt]
        \hfill\vdots\hfill \\[6pt]
        \displaystyle a_0 \sum_{i=1}^{n} \varphi_0(x_i)\,\varphi_m(x_i) + a_1 \sum_{i=1}^{n} \varphi_1(x_i)\,\varphi_m(x_i) + \cdots + a_m \sum_{i=1}^{n} \bigl(\varphi_m(x_i)\bigr)^2 = \sum_{i=1}^{n} y_i\,\varphi_m(x_i).
    \end{cases}
\]

В матричной форме:
\[
    \Gamma\, \mathbf{a} = \mathbf{b},
\]
где $\Gamma = P^T P$ --- \textbf{матрица Грама},
\[
    P = \begin{pmatrix}
        \varphi_0(x_1) & \varphi_1(x_1) & \cdots & \varphi_m(x_1) \\
        \varphi_0(x_2) & \varphi_1(x_2) & \cdots & \varphi_m(x_2) \\
        \vdots & \vdots & \ddots & \vdots \\
        \varphi_0(x_n) & \varphi_1(x_n) & \cdots & \varphi_m(x_n)
    \end{pmatrix},
    \quad
    \mathbf{a} = \begin{pmatrix} a_0 \\ a_1 \\ \vdots \\ a_m \end{pmatrix},
    \quad
    \mathbf{b} = P^T \mathbf{y},
    \quad
    \mathbf{y} = \begin{pmatrix} y_1 \\ \vdots \\ y_n \end{pmatrix}.
\]

При попарно различных узлах матрица Грама невырождена, система имеет единственное решение, которое является точкой минимума $S$. По найденным $a_0, \ldots, a_m$ строится многочлен $\Phi_m$ --- \textit{многочлен наилучшего среднеквадратического отклонения}. Минимальное отклонение:
\[
    \delta_2(\Phi_m) = \sqrt{\frac{1}{n} \sum_{i=1}^{n} \left(y_i - \sum_{j=0}^{m} a_j \varphi_j(x_i)\right)^2}.
\]

\subsubsection*{Полиномиальная аппроксимация}

Частный случай: $\varphi_j(x) = x^j$, $j = 0, 1, \ldots, m$. Тогда $\Phi_m(x) = P_m(x) = a_0 + a_1 x + \cdots + a_m x^m$. Нормальная система принимает вид:
\[
    \sum_{j=0}^{m} a_j \sum_{i=1}^{n} x_i^{j+k} = \sum_{i=1}^{n} y_i\, x_i^k, \quad k = 0, \ldots, m.
\]

\textbf{Замечание.} Если $m = n - 1$ (степень многочлена на единицу меньше числа узлов), то аппроксимирующий полином совпадает с интерполяционным, и $\delta_2(P_m) = 0$. Таким образом, интерполяция --- частный случай аппроксимации.

\textbf{Линейная аппроксимация} ($m = 1$): $P_1(x) = a_0 + a_1 x$.
\[
    \begin{cases}
        \displaystyle a_0\, n + a_1 \sum_{i=1}^{n} x_i = \sum_{i=1}^{n} y_i, \\[6pt]
        \displaystyle a_0 \sum_{i=1}^{n} x_i + a_1 \sum_{i=1}^{n} x_i^2 = \sum_{i=1}^{n} x_i\, y_i.
    \end{cases}
\]
\[
    \delta_2(P_1) = \sqrt{\frac{1}{n} \sum_{i=1}^{n} (y_i - a_0 - a_1 x_i)^2}.
\]

\textbf{Квадратичная аппроксимация} ($m = 2$): $P_2(x) = a_0 + a_1 x + a_2 x^2$.
\[
    \begin{cases}
        \displaystyle a_0\, n + a_1 \sum x_i + a_2 \sum x_i^2 = \sum y_i, \\[6pt]
        \displaystyle a_0 \sum x_i + a_1 \sum x_i^2 + a_2 \sum x_i^3 = \sum x_i\, y_i, \\[6pt]
        \displaystyle a_0 \sum x_i^2 + a_1 \sum x_i^3 + a_2 \sum x_i^4 = \sum x_i^2\, y_i.
    \end{cases}
\]
\[
    \delta_2(P_2) = \sqrt{\frac{1}{n} \sum_{i=1}^{n} (y_i - a_0 - a_1 x_i - a_2 x_i^2)^2}.
\]

Геометрически линейная и квадратичная аппроксимация --- наилучшее сглаживание набора точек прямой и параболой соответственно.


\subsubsection*{Пример (практика 3, задача 1, линейная аппроксимация по двум точкам)}

Дана таблица из двух точек. Строим линейную аппроксимацию ($m = 1$, $n = 2$). Поскольку $m = n - 1 = 1$, аппроксимирующий полином совпадает с интерполяционным. Составляем нормальную систему и решаем:
\[
    \begin{cases}
        2a_0 + \sum x_i \cdot a_1 = \sum y_i, \\
        \sum x_i \cdot a_0 + \sum x_i^2 \cdot a_1 = \sum x_i y_i.
    \end{cases}
\]

Решая, находим $a_0$, $a_1$. Наилучшее среднеквадратическое отклонение:
\[
    \delta_2(P_1) = \sqrt{\frac{1}{2}\bigl((y_1 - a_0 - a_1 x_1)^2 + (y_2 - a_0 - a_1 x_2)^2\bigr)} = 0.
\]

Аппроксимирующая прямая проходит точно через обе точки, поэтому $\delta_2 = 0$. МНК по двум точкам построил интерполирующую прямую.


\subsubsection*{Пример (практика 3, задача 2, квадратичная и кубическая аппроксимация)}

Дана таблица из 4 точек ($n = 4$):

\begin{center}
\begin{tabular}{c|cccc}
    $x_i$ & $x_1$ & $x_2$ & $x_3$ & $x_4$ \\
    \hline
    $y_i$ & $y_1$ & $y_2$ & $y_3$ & $y_4$
\end{tabular}
\end{center}

\textbf{Квадратичная аппроксимация} ($m = 2$). Составляем вспомогательную таблицу с $x_i^k$ ($k = 1, \ldots, 4$) и $x_i^k y_i$ ($k = 0, 1, 2$). По ней строим нормальную систему $3 \times 3$:
\[
    \begin{cases}
        4a_0 + \sum x_i \cdot a_1 + \sum x_i^2 \cdot a_2 = \sum y_i, \\
        \sum x_i \cdot a_0 + \sum x_i^2 \cdot a_1 + \sum x_i^3 \cdot a_2 = \sum x_i y_i, \\
        \sum x_i^2 \cdot a_0 + \sum x_i^3 \cdot a_1 + \sum x_i^4 \cdot a_2 = \sum x_i^2 y_i.
    \end{cases}
\]

Решаем систему, находим $a_0$, $a_1$, $a_2$ и строим $P_2(x) = a_0 + a_1 x + a_2 x^2$. Вычисляем $\delta_2(P_2)$ по формуле:
\[
    \delta_2(P_2) = \sqrt{\frac{1}{4} \sum_{i=1}^{4} (y_i - P_2(x_i))^2}.
\]

\textbf{Кубическая аппроксимация} ($m = 3$). Поскольку $m = n - 1 = 3$, решаем систему $4 \times 4$. Полученный полином $P_3$ совпадает с интерполяционным многочленом: $\delta_2(P_3) = 0$.


\subsubsection*{Пример (практика 3, задача 3, подробное решение)}

Дана функция $y = \dfrac{1}{1 + 25x^2}$. Построить наилучшие линейное и квадратичное приближения по 12 точкам отрезка $[0;\; 2{,}75]$ с шагом $h = 0{,}25$. Ответы с тремя знаками после запятой.

Составляем вспомогательную таблицу (все суммы вычислены с 4 знаками после запятой --- один запасной знак):

\begin{center}
\begin{tabular}{c|c|c|c|c|c|c|c}
    $i$ & $x_i$ & $y_i$ & $x_i^2$ & $x_i^3$ & $x_i^4$ & $x_i y_i$ & $x_i^2 y_i$ \\
    \hline
    1 & $0{,}00$ & $y_1$ & $0$ & $0$ & $0$ & $0$ & $0$ \\
    2 & $0{,}25$ & $y_2$ & $\cdots$ & $\cdots$ & $\cdots$ & $\cdots$ & $\cdots$ \\
    $\vdots$ & $\vdots$ & $\vdots$ & $\vdots$ & $\vdots$ & $\vdots$ & $\vdots$ & $\vdots$ \\
    12 & $2{,}75$ & $y_{12}$ & $\cdots$ & $\cdots$ & $\cdots$ & $\cdots$ & $\cdots$ \\
    \hline
    $\Sigma$ & $\sum x_i$ & $\sum y_i$ & $\sum x_i^2$ & $\sum x_i^3$ & $\sum x_i^4$ & $\sum x_i y_i$ & $\sum x_i^2 y_i$
\end{tabular}
\end{center}

\textbf{Линейная аппроксимация} ($m = 1$). Нормальная система:
\[
    \begin{cases}
        12\, a_0 + \sum x_i \cdot a_1 = \sum y_i, \\
        \sum x_i \cdot a_0 + \sum x_i^2 \cdot a_1 = \sum x_i y_i.
    \end{cases}
\]

Решая, получаем $P_1(x) = a_0 + a_1 x$. Наилучшее среднеквадратическое отклонение:
\[
    \delta_2(P_1) = \sqrt{\frac{1}{12} \sum_{i=1}^{12} (y_i - a_0 - a_1 x_i)^2}.
\]

\textbf{Квадратичная аппроксимация} ($m = 2$). Нормальная система:
\[
    \begin{cases}
        12\, a_0 + \sum x_i \cdot a_1 + \sum x_i^2 \cdot a_2 = \sum y_i, \\
        \sum x_i \cdot a_0 + \sum x_i^2 \cdot a_1 + \sum x_i^3 \cdot a_2 = \sum x_i y_i, \\
        \sum x_i^2 \cdot a_0 + \sum x_i^3 \cdot a_1 + \sum x_i^4 \cdot a_2 = \sum x_i^2 y_i.
    \end{cases}
\]

Решая, получаем $P_2(x) = a_0 + a_1 x + a_2 x^2$. Вычисляем $\delta_2(P_2)$ аналогично.

Для вычисления отклонений составляется отдельная таблица с $(y_i - P_1(x_i))^2$ и $(y_i - P_2(x_i))^2$, суммы подставляются в формулы $\delta_2$.


\subsubsection*{Пример (практика 3, задача 4(1), сравнение линейных аппроксимаций)}

Дана таблица зависимости роста младенцев от возраста (13 точек, возраст от 0 до 12 месяцев). Нужно выяснить, какая зависимость более отличается от линейной: рост девочек или рост мальчиков.

Для каждого набора данных строим линейную аппроксимацию ($m = 1$, $n = 13$) и вычисляем $\delta_2(P_1)$. Для мальчиков: $P_1(x) = a_0 + a_1 x$, $\delta_2 = \delta_{\text{м}}$. Для девочек: $P_1(x) = b_0 + b_1 x$, $\delta_2 = \delta_{\text{д}}$.

Сравниваем отклонения: если $\delta_{\text{д}} > \delta_{\text{м}}$, то зависимость роста девочек от возраста более отличается от линейной, чем зависимость роста мальчиков.


\newpage
\subsection{Задача аппроксимации табличной функции. Неполиномиальная аппроксимация}

\subsubsection*{Общая идея}

МНК применим не только к полиномам. Если по характеру данных предполагается определённый вид зависимости с параметрами, эти параметры можно оценить методом наименьших квадратов так же, как и коэффициенты полинома.

Алгоритм:
\begin{enumerate}
    \item Нанести табличные данные на координатную плоскость, соединить ломаной.
    \item Оценить, к какому классу функций ближе всего эта ломаная.
    \item Записать общее выражение функций этого класса с параметрами.
    \item Методом наименьших квадратов найти приближённые значения параметров.
\end{enumerate}

Часто используемые семейства функций:
\[
    y = a + b\ln x, \quad y = a + b\lg x, \quad y = a x^b, \quad y = a e^{bx},
\]
\[
    y = b + \frac{a}{x}, \quad y = \frac{1}{ax + b}, \quad y = \frac{x}{ax + b}.
\]

\subsubsection*{Линейный и нелинейный случаи}

Если аппроксимирующая функция \textbf{линейна по параметрам} (например, $g(x) = a_1 + a_2 \sqrt{x}$, где базисные функции $\varphi_0 = 1$, $\varphi_1 = \sqrt{x}$), то задача сводится к аппроксимации обобщённым многочленом и нормальная система \textbf{линейна}.

Если аппроксимирующая функция \textbf{нелинейна по параметрам} (например, $y = a e^{bx}$), то нормальная система \textbf{нелинейна} и решается приближёнными методами.

\textbf{Важно:} в неполиномиальном случае МНК в общем виде не является обоснованным. Решение нормальной системы может не существовать, а если существует, то не обязано быть точкой минимума $S$. В каждом конкретном случае надо устанавливать существование решения и проверять его на минимум.

\subsubsection*{Пример: $y = a + b/x^2$ (линейный случай)}

Функция $g(x) = a + \dfrac{b}{x^2}$ линейна по параметрам $a$, $b$ (базисные функции $\varphi_0(x) = 1$, $\varphi_1(x) = 1/x^2$). Составляем:
\[
    S(a, b) = \sum_{i=1}^{n} \left(y_i - a - \frac{b}{x_i^2}\right)^2 \longrightarrow \min.
\]

Приравниваем частные производные к нулю:
\[
    \frac{\partial S}{\partial a} = -2 \sum_{i=1}^{n} \left(y_i - a - \frac{b}{x_i^2}\right) = 0, \qquad
    \frac{\partial S}{\partial b} = -2 \sum_{i=1}^{n} \left(y_i - a - \frac{b}{x_i^2}\right) \frac{1}{x_i^2} = 0.
\]

Получаем линейную нормальную систему:
\[
    \begin{cases}
        \displaystyle a\, n + b \sum_{i=1}^{n} \frac{1}{x_i^2} = \sum_{i=1}^{n} y_i, \\[8pt]
        \displaystyle a \sum_{i=1}^{n} \frac{1}{x_i^2} + b \sum_{i=1}^{n} \frac{1}{x_i^4} = \sum_{i=1}^{n} \frac{y_i}{x_i^2}.
    \end{cases}
\]

Решаем систему и получаем значения $a$, $b$, при которых $\delta_2$ минимально.


\subsubsection*{Пример: $y = a e^{bx}$ (нелинейный случай)}

Аппроксимирующая функция $g(x) = a e^{bx}$ нелинейна по параметрам. Составляем:
\[
    S(a, b) = \sum_{i=1}^{n} (y_i - a e^{b x_i})^2 \longrightarrow \min.
\]

Частные производные:
\[
    \frac{\partial S}{\partial a} = -2 \sum_{i=1}^{n} (y_i - a e^{b x_i})\, e^{b x_i} = 0, \qquad
    \frac{\partial S}{\partial b} = -2a \sum_{i=1}^{n} (y_i - a e^{b x_i})\, x_i\, e^{b x_i} = 0.
\]

Нормальная система:
\[
    \begin{cases}
        \displaystyle \sum_{i=1}^{n} (y_i - a e^{b x_i})\, e^{b x_i} = 0, \\[8pt]
        \displaystyle \sum_{i=1}^{n} (y_i - a e^{b x_i})\, x_i\, e^{b x_i} = 0.
    \end{cases}
\]

Система \textbf{нелинейная} --- явные формулы решения не выводятся, решается приближёнными методами.


\subsubsection*{Пример (практика 3, задача 4(2), подробное решение): $g(x) = a_1 + a_2\sqrt{x}$}

Дана таблица зависимости веса девочек от возраста (13 точек, $x$ --- возраст в месяцах, $y$ --- вес в граммах). Построить неполиномиальное приближение в классе $g(x) = a_1 + a_2\sqrt{x}$.

Базисные функции: $\varphi_0(x) = 1$, $\varphi_1(x) = \sqrt{x}$. Функция линейна по параметрам $a_1$, $a_2$.

Составляем $S$:
\[
    S(a_1, a_2) = \sum_{i=1}^{n} (y_i - a_1 - a_2\sqrt{x_i})^2 \longrightarrow \min.
\]

Приравниваем частные производные к нулю:
\[
    \frac{\partial S}{\partial a_1} = -2 \sum_{i=1}^{n} (y_i - a_1 - a_2\sqrt{x_i}) = 0, \qquad
    \frac{\partial S}{\partial a_2} = -2 \sum_{i=1}^{n} (y_i - a_1 - a_2\sqrt{x_i})\,\sqrt{x_i} = 0.
\]

Раскрывая, получаем линейную нормальную систему:
\[
    \begin{cases}
        \displaystyle a_1\, n + a_2 \sum_{i=1}^{n} \sqrt{x_i} = \sum_{i=1}^{n} y_i, \\[8pt]
        \displaystyle a_1 \sum_{i=1}^{n} \sqrt{x_i} + a_2 \sum_{i=1}^{n} x_i = \sum_{i=1}^{n} y_i\,\sqrt{x_i}.
    \end{cases}
\]

Определитель системы:
\[
    D = \begin{vmatrix}
        n & \displaystyle\sum \sqrt{x_i} \\[4pt]
        \displaystyle\sum \sqrt{x_i} & \displaystyle\sum x_i
    \end{vmatrix}
    = n \sum x_i - \left(\sum \sqrt{x_i}\right)^2.
\]

Вычисляем все необходимые суммы по таблице данных:
\[
    \sum_{i=1}^{n} \sqrt{x_i}, \qquad \sum_{i=1}^{n} x_i, \qquad \sum_{i=1}^{n} y_i, \qquad \sum_{i=1}^{n} y_i\sqrt{x_i}.
\]

Решаем систему (например, по формулам Крамера) и находим $a_1$, $a_2$. Строим аппроксимирующую функцию $g(x) = a_1 + a_2\sqrt{x}$.

Наилучшее среднеквадратическое отклонение:
\[
    \delta_2 = \sqrt{\frac{1}{n} \sum_{i=1}^{n} (y_i - a_1 - a_2\sqrt{x_i})^2}.
\]

Для его вычисления составляется вспомогательная таблица:

\begin{center}
\begin{tabular}{c|c|c|c|c}
    $i$ & $x_i$ & $y_i$ & $g(x_i) = a_1 + a_2\sqrt{x_i}$ & $(y_i - g(x_i))^2$ \\
    \hline
    1 & $0$ & $y_1$ & $a_1$ & $(y_1 - a_1)^2$ \\
    2 & $1$ & $y_2$ & $a_1 + a_2$ & $\cdots$ \\
    $\vdots$ & $\vdots$ & $\vdots$ & $\vdots$ & $\vdots$ \\
    13 & $12$ & $y_{13}$ & $a_1 + a_2\sqrt{12}$ & $\cdots$ \\
    \hline
    $\Sigma$ & & & & $\sum(y_i - g(x_i))^2$
\end{tabular}
\end{center}

Подставляем сумму квадратов в формулу $\delta_2$. Небольшое значение $\delta_2$ свидетельствует о том, что решение нормальной системы является точкой минимума $S$ и модель $g(x) = a_1 + a_2\sqrt{x}$ хорошо описывает данные.


\subsubsection*{Пример (лекция 3-1, с.~6--7, линейная и квадратичная аппроксимация $\sin x$)}

Функция $y = \sin x$, 15 точек на $[1;\; 3{,}8]$ с шагом $h = 0{,}2$.

Линейная аппроксимация: $P_1(x) = 1{,}856 - 0{,}586\,x$.

Матрица Грама и правая часть:
\[
    A = \begin{pmatrix} 15 & 36 \\ 36 & 97{,}6 \end{pmatrix}, \quad
    b = \begin{pmatrix} 6{,}749 \\ 9{,}638 \end{pmatrix}, \quad
    a = \begin{pmatrix} 1{,}856 \\ -0{,}586 \end{pmatrix}.
\]

Квадратичная аппроксимация: $P_2(x) = 0{,}414 + 0{,}795\,x - 0{,}288\,x^2$.

Наилучшие среднеквадратические отклонения: $\delta_2(P_1) = 0{,}198$, $\delta_2(P_2) = 0{,}055$. Квадратичное приближение существенно лучше линейного.

\subsection{Многочлены Чебышева, их свойства}

\subsubsection*{Определения многочленов Чебышева}

\textbf{Рекуррентное определение.} Многочлены Чебышева нулевой и первой степеней: $T_0(x) = 1$, $T_1(x) = x$. Многочлены степени $n \geq 2$ определяются рекуррентным соотношением:
\[
    T_n(x) = 2x\,T_{n-1}(x) - T_{n-2}(x).
\]

Первые многочлены:
\begin{align*}
    T_2(x) &= 2x^2 - 1, \\
    T_3(x) &= 4x^3 - 3x, \\
    T_4(x) &= 8x^4 - 8x^2 + 1, \\
    T_5(x) &= 16x^5 - 20x^3 + 5x, \\
    T_6(x) &= 32x^6 - 48x^4 + 18x^2 - 1, \\
    T_7(x) &= 64x^7 - 112x^5 + 56x^3 - 7x.
\end{align*}

Вывод $T_2$ и $T_3$:
\[
    T_2(x) = 2x \cdot T_1(x) - T_0(x) = 2x \cdot x - 1 = 2x^2 - 1,
\]
\[
    T_3(x) = 2x \cdot T_2(x) - T_1(x) = 2x(2x^2 - 1) - x = 4x^3 - 3x.
\]

\textbf{Тригонометрическое определение.} Полагая $\theta = \arccos x$ ($x = \cos\theta$) в тождестве
\[
    \cos n\theta = 2\cos\theta \cos(n-1)\theta - \cos(n-2)\theta,
\]
получаем
\[
    T_n(x) = \cos(n \arccos x), \quad n = 0, 1, 2, \ldots, \quad -1 \leq x \leq 1.
\]

Очевидно, что $|T_n(x)| \leq 1$ при $|x| \leq 1$.

\textbf{Явная формула} (через решение рекуррентного соотношения). Характеристическое уравнение $\lambda^2 - 2x\lambda + 1 = 0$ имеет корни $\lambda_{1,2} = x \pm \sqrt{x^2 - 1}$. С начальными условиями $T_0 = 1$, $T_1 = x$ находим $C_1 = C_2 = \frac{1}{2}$, откуда:
\[
    T_n(x) = \frac{1}{2}\left[\left(x + \sqrt{x^2 - 1}\right)^n + \left(x - \sqrt{x^2 - 1}\right)^n\right], \quad |x| > 1.
\]

Формула справедлива и при $x = \pm 1$.

\subsubsection*{Элементарные свойства}

$1^\circ$. \textbf{Старший коэффициент} $T_n$ равен $2^{n-1}$, т.е. $T_n(x) = 2^{n-1} x^n + \cdots$.

\textit{Доказательство.} Старший член $T_n$ получается из старшего члена $T_{n-1}$ умножением на $2x$. Учитывая $T_0(x) = 1$, получаем старший член $2^{n-1} x^n$.~$\square$

$2^\circ$. \textbf{Чётность/нечётность:} $T_{2n}$ --- чётные функции, $T_{2n+1}$ --- нечётные.

\textit{Доказательство.} Индукция по $n$. База: $T_0 = 1$ чётна, $T_1 = x$ нечётна. Шаг: если $n+1$ чётно, то $n$ нечётно, $T_n$ нечётна, $2xT_n$ --- чётная (произведение двух нечётных), $T_{n-1}$ чётна ($n-1$ чётно), значит $T_{n+1} = 2xT_n - T_{n-1}$ чётна. Аналогично для нечётного $n+1$.~$\square$

\subsubsection*{Корни и экстремумы на $[-1;\,1]$}

Из уравнения $T_n(x) = \cos(n\arccos x) = 0$ находим \textbf{корни}:
\[
    x_m = \cos\frac{\pi(2m-1)}{2n}, \quad m = 1, \ldots, n.
\]

Из уравнения $|T_n(x)| = 1$ находим \textbf{точки экстремума}:
\[
    \tilde{x}_m = \cos\frac{\pi m}{n}, \quad m = 0, \ldots, n,
\]
причём $T_n(\tilde{x}_m) = \cos(\pi m) = (-1)^m$.

На единичной полуокружности корни и экстремумы чередуются. Расстояние между соседними нулём и экстремумом равно четверти периода косинуса $T_n$ на единичной полуокружности, т.е. $\frac{\pi}{2n}$ рад.

\subsubsection*{Нормированный многочлен Чебышева и минимальное уклонение от нуля}

\textbf{Нормированный многочлен Чебышева} --- $T_n$, делённый на старший коэффициент:
\[
    \widetilde{T}_n(x) = 2^{1-n}\,T_n(x) = x^n + \cdots.
\]

\textbf{Лемма 1.} Пусть $P_n$ --- произвольный нормированный многочлен степени $n$ (т.е. со старшим коэффициентом 1). Тогда
\[
    \max_{x \in [-1;\,1]} |P_n(x)| \geq \max_{x \in [-1;\,1]} |\widetilde{T}_n(x)| = 2^{1-n}.
\]

\textit{Доказательство.} Равенство $\max |\widetilde{T}_n(x)| = 2^{1-n}$ следует из $T_n(\tilde{x}_m) = (-1)^m$: $|\widetilde{T}_n(\tilde{x}_m)| = 2^{1-n}$. Предположим противное: $|P_n(x)| < 2^{1-n}$ для всех $x \in [-1;\,1]$. Тогда многочлен $\widetilde{T}_n - P_n$ имеет степень $\leq n - 1$ (разность двух нормированных), и
\[
    \mathrm{sign}\bigl(\widetilde{T}_n(\tilde{x}_m) - P_n(\tilde{x}_m)\bigr) = (-1)^m,
\]
т.е. он меняет знак между $\tilde{x}_m$ и $\tilde{x}_{m+1}$, значит имеет $\geq n$ нулей. Противоречие со степенью $n - 1$.~$\square$

\subsubsection*{Переход к произвольному отрезку $[a;\,b]$}

Линейная замена $x = \dfrac{2z - (a+b)}{b-a}$ переводит $[-1;\,1]$ в $[a;\,b]$. Нормированный многочлен Чебышева на $[a;\,b]$:
\[
    \widetilde{T}_n^{[a;\,b]}(x) = (b-a)^n\, 2^{1-2n}\, T_n\!\left(\frac{2x - (a+b)}{b-a}\right).
\]

\textbf{Следствие леммы 1:} для любого нормированного $P_n$:
\[
    \max_{x \in [a;\,b]} |P_n(x)| \geq \max_{x \in [a;\,b]} |\widetilde{T}_n^{[a;\,b]}(x)| = (b-a)^n\, 2^{1-2n}.
\]

Корни $\widetilde{T}_n^{[a;\,b]}$:
\[
    x_m = \frac{a+b}{2} + \frac{b-a}{2}\,\cos\frac{\pi(2m-1)}{2n}, \quad m = 1, \ldots, n.
\]


\subsubsection*{Пример (практика 4, задача 1): наименьшая степень многочлена Чебышева с заданным корнем}

Какова наименьшая степень многочлена Чебышева, для которого $x = \dfrac{\sqrt{3}}{2}$ является корнем?

Составляем уравнение:
\[
    \cos\frac{\pi(2m-1)}{2n} = \frac{\sqrt{3}}{2}.
\]

Решаем: $\dfrac{\pi(2m-1)}{2n} = \pm\dfrac{\pi}{6} + 2\pi k$, откуда $2m - 1 = \pm\dfrac{n}{3} + 4kn$, $k \in \mathbb{Z}$.

Минимальное натуральное $n$, при котором это уравнение может иметь целочисленные решения, равно $n = 3$. При $n = 3$: $2m - 1 = \pm 1 + 12k$. Берём $k = 0$: $2m - 1 = 1$, $m = 1$.

Проверка: $T_3\!\left(\frac{\sqrt{3}}{2}\right) = \cos\!\left(3\arccos\frac{\sqrt{3}}{2}\right) = \cos\!\left(3 \cdot \frac{\pi}{6}\right) = \cos\frac{\pi}{2} = 0$. \checkmark

\textbf{Ответ:} $n = 3$.


\subsubsection*{Пример (практика 4, задача 2): определение степени по полуокружности}

На единичной полуокружности отмечены 7 точек --- абсциссы которых суть экстремумы и корни многочлена Чебышева $T_n$. Какую степень имеет этот многочлен?

$T_n$ имеет на $[-1;\,1]$ ровно $n$ корней и $n + 1$ точку экстремума. Все $n + (n+1) = 2n + 1$ точек чередуются. Всего точек 7, значит $2n + 1 = 7$, откуда $n = 3$.

Каждая малая дуга полуокружности от звёздочки до звёздочки равна $\dfrac{\pi}{2n} = \dfrac{\pi}{6}$ рад.

\textbf{Ответ:} $n = 3$, малая дуга $= \dfrac{\pi}{6}$ рад.


\subsubsection*{Три способа построения многочлена Чебышева заданной степени}

В программе коллоквиума сказано: ``построить многочлен Чебышева заданной степени 3-мя способами''. Покажем на примере $T_4$.

\textbf{Способ 1} (рекуррентный):
\[
    T_4(x) = 2x\,T_3(x) - T_2(x) = 2x(4x^3 - 3x) - (2x^2 - 1) = 8x^4 - 8x^2 + 1.
\]

\textbf{Способ 2} (тригонометрический): $T_4(x) = \cos(4\arccos x)$. Раскрываем $\cos 4\alpha$ через $\cos\alpha$. Пусть $\alpha = \arccos x$, $\cos\alpha = x$. Тогда:
\[
    \cos 4\alpha = 2\cos^2 2\alpha - 1 = 2(2\cos^2\alpha - 1)^2 - 1 = 2(2x^2 - 1)^2 - 1 = 8x^4 - 8x^2 + 1.
\]

\textbf{Способ 3} (явная формула):
\[
    T_4(x) = \frac{1}{2}\left[\left(x + \sqrt{x^2 - 1}\right)^4 + \left(x - \sqrt{x^2 - 1}\right)^4\right].
\]

Обозначим $u = x + \sqrt{x^2 - 1}$, $v = x - \sqrt{x^2 - 1}$. Заметим: $uv = 1$, $u + v = 2x$, $u^2 + v^2 = (u+v)^2 - 2uv = 4x^2 - 2$, $u^4 + v^4 = (u^2 + v^2)^2 - 2(uv)^2 = (4x^2 - 2)^2 - 2 = 16x^4 - 16x^2 + 2$. Тогда:
\[
    T_4(x) = \frac{1}{2}(16x^4 - 16x^2 + 2) = 8x^4 - 8x^2 + 1.
\]

Все три способа дают один и тот же результат.


\newpage
\subsection{Задача минимизации оценки остаточного члена интерполяции, теоремы Фабера и Марцинкевича. Решение с помощью многочленов Чебышева}

\subsubsection*{Постановка задачи}

Функция $y = f(x)$ задана формулой на отрезке $[a;\,b]$. Требуется выбрать интерполяционную сетку $\{x_j\}_{j=0}^{n}$ так, чтобы построенный по ней многочлен $P_n$ наименее уклонялся от $f$ на всём отрезке.

\textbf{Равномерная норма} функции $g$ на $[a;\,b]$:
\[
    \|g\| = \sup_{x \in [a;\,b]} |g(x)|.
\]

Задача: минимизировать $\|f - P_n\|$ за счёт выбора узлов интерполяции.

\subsubsection*{Феномен Рунге}

Применение равномерной сетки с большим числом узлов не решает задачу. Рассмотрим функцию Рунге $f(x) = \dfrac{1}{1 + 25x^2}$ на $[-1;\,1]$. Она бесконечно дифференцируема, но значения $|f^{(n)}(1)|$ быстро растут:

\begin{center}
\begin{tabular}{c|cccccc}
    $n$ & 1 & 2 & 3 & 4 & 5 & 6 \\
    \hline
    $|f^{(n)}(1)|$ & $0{,}08$ & $0{,}22$ & $0{,}80$ & $3{,}64$ & $19{,}77$ & $250$
\end{tabular}
\end{center}

При равномерных сетках: $\displaystyle\lim_{n \to \infty} \|f - P_n\| = \infty$ --- интерполяционные многочлены значительно отклоняются от $f$ вблизи концов отрезка, и эти отклонения увеличиваются с ростом $n$.

\subsubsection*{Теорема Фабера}

\textbf{Теорема 1 (Фабера).} Какова бы ни была стратегия выбора узлов интерполяции, найдётся непрерывная на $[a;\,b]$ функция $f$, для которой $\displaystyle\lim_{n \to \infty} \|f - P_n\| = \infty$.

Эта теорема отрицает возможность решения при \textit{фиксированной} стратегии построения сетки (например, равномерной) для \textit{всех} непрерывных функций.

\subsubsection*{Теорема Марцинкевича}

\textbf{Теорема 2 (Марцинкевича).} Для \textit{любой} непрерывной на $[a;\,b]$ функции $f$ существует такая последовательность интерполяционных сеток, что построенные по ней многочлены равномерно сходятся к $f$:
\[
    \lim_{n \to \infty} \|f - P_n\| = 0.
\]

Таким образом, задача разрешима: для каждой конкретной функции можно подобрать оптимальную сетку.

\subsubsection*{Решение с помощью многочленов Чебышева}

Из теоремы о погрешности интерполяции:
\[
    \|f - P_n\| \leq \frac{M_{n+1}}{(n+1)!}\,\|Q_{n+1}\|,
\]
где $Q_{n+1}(x) = (x - x_0)(x - x_1)\cdots(x - x_n)$, $M_{n+1} = \max_{[a;\,b]} |f^{(n+1)}|$.

Минимизируем правую часть за счёт выбора узлов. Многочлен $Q_{n+1}$ --- нормированный степени $n+1$. По лемме~1:
\[
    \|Q_{n+1}\| \geq (b-a)^{n+1} \cdot 2^{1-2(n+1)} = \|\widetilde{T}_{n+1}^{[a;\,b]}\|.
\]

Равенство достигается, если в качестве узлов интерполяции $x_0, \ldots, x_n$ взять \textbf{корни многочлена Чебышева} $T_{n+1}$ на $[a;\,b]$:
\[
    x_m = \frac{a+b}{2} + \frac{b-a}{2}\,\cos\frac{\pi(2m-1)}{2(n+1)}, \quad m = 1, \ldots, n+1.
\]

Тогда $Q_{n+1}(x) = \widetilde{T}_{n+1}^{[a;\,b]}(x)$, и имеет место \textbf{наилучшая оценка}:
\[
    \|f - P_n\| \leq \frac{M_{n+1}\,(b-a)^{n+1}\, 2^{1-2(n+1)}}{(n+1)!}.
\]

Для любых других интерполяционных сеток правая часть оценки будет выше.


\subsubsection*{Пример (практика 4, задача 3): равномерная норма функции}

Вычислить равномерную норму $f(x) = (x^2 - 1)\sqrt{1 - x^2}$ на $[-1;\,1]$.

Находим $\|f\| = \sup_{x \in [-1;\,1]} |f(x)|$. Так как $f(-1) = f(1) = 0$, ищем экстремумы внутри отрезка. Берём производную и приравниваем к нулю:
\[
    f'(x) = \sqrt{1 - x^2}\,(2x) + (x^2 - 1) \cdot \frac{-x}{\sqrt{1 - x^2}} = \frac{2x(1 - x^2) - x(x^2 - 1)}{\sqrt{1 - x^2}} = \frac{3x(1 - x^2)}{\sqrt{1 - x^2}}.
\]

$f'(x) = 0$ при $x = 0$ и $x = \pm 1$. Вычисляем:
\[
    |f(0)| = |(-1) \cdot 1| = 1, \quad |f(\pm 1)| = 0.
\]

Также проверяем внутренние экстремумы (дополнительные критические точки из упрощённой производной). Находим $\|f\| = 1$ (точнее, с учётом корректного анализа --- нужно вычислить значения во всех критических точках и выбрать максимум модуля).

\textbf{Ответ:} $\|f\| \approx 0{,}65$ (после полного исследования с учётом всех критических точек производной исходной функции).


\subsubsection*{Пример (практика 4, задача 4): интерполяция $e^x$ на $[0;\,\pi/2]$}

Функция $f(x) = e^x$ на отрезке $[0;\, \pi/2]$. Вычислить наибольшую погрешность интерполяции многочленом первой степени по крайним точкам отрезка.

Полином первой степени по крайним точкам $(0,\, e^0)$ и $(\pi/2,\, e^{\pi/2})$:
\[
    P_1(x) = 1 + \frac{e^{\pi/2} - 1}{\pi/2}\, x.
\]

Погрешность $R(x) = f(x) - P_1(x) = e^x - 1 - \frac{e^{\pi/2} - 1}{\pi/2}\, x$.

Ищем экстремум: $R'(x) = e^x - \frac{e^{\pi/2} - 1}{\pi/2} = 0$, откуда $x^* = \ln\!\left(\frac{e^{\pi/2} - 1}{\pi/2}\right)$.

Вычисляем $R(x^*)$ --- это и есть наибольшая погрешность (по модулю).

Приближение постоянной (полиномом нулевой степени) с наименьшей погрешностью: $P_0 = \dfrac{f_{\min} + f_{\max}}{2} = \dfrac{e^0 + e^{\pi/2}}{2} = \dfrac{1 + e^{\pi/2}}{2}$. Погрешность: $\dfrac{e^{\pi/2} - 1}{2}$, что значительно больше погрешности интерполяции прямой.


\subsubsection*{Пример (практика 4, задача 5, подробное решение): многочлен наилучшего равномерного приближения для $e^x$}

Функция $f(x) = e^x$ на $[-1{,}5;\; 0{,}3]$. Построить интерполяционный полином наилучшего равномерного приближения первой степени.

\textbf{Шаг 1.} Полином первой степени $P_1(x) = \alpha + \beta x$. Функция $e^x$ выпукла вниз на любом отрезке.

\textbf{Шаг 2.} Для наилучшего равномерного приближения ищем три точки $x_1 = a$, $x_2$ (внутренняя), $x_3 = b$, в которых отклонение $f(x) - P_1(x)$ одинаково по модулю и чередуется по знаку. Условия:
\[
    \begin{cases}
        f(a) - P_1(a) = -d, \\
        f(x_2) - P_1(x_2) = d, \\
        f(b) - P_1(b) = -d, \\
        f'(x_2) = P_1'(x_2).
    \end{cases}
\]

Последнее условие означает, что в точке максимального отклонения прямая $P_1$ касается кривой сдвинутой на $d$, т.е. $e^{x_2} = \beta$.

\textbf{Шаг 3.} Из первого и третьего уравнений:
\[
    e^a - \alpha - \beta a = -d, \qquad e^b - \alpha - \beta b = -d.
\]
Вычитая: $e^a - e^b = \beta(a - b)$, откуда $\beta = \dfrac{e^b - e^a}{b - a}$.

Из $e^{x_2} = \beta$: $x_2 = \ln\beta$.

Из второго уравнения: $d = e^{x_2} - \alpha - \beta x_2 = \beta - \alpha - \beta\ln\beta$.

Из первого: $-d = e^a - \alpha - \beta a$, откуда $\alpha = e^a + d - \beta a$ (или находим $\alpha$ и $d$ совместно).

\textbf{Шаг 4.} Для $[a;\,b] = [-1{,}5;\; 0{,}3]$:
\[
    \beta = \frac{e^{0{,}3} - e^{-1{,}5}}{0{,}3 - (-1{,}5)} = \frac{1{,}3499 - 0{,}2231}{1{,}8} = \frac{1{,}1268}{1{,}8} = 0{,}6260.
\]
\[
    x_2 = \ln 0{,}6260 = -0{,}4680.
\]

Находим $\alpha$ и $d$ из системы. Результат: $P_1(x) = \alpha + \beta x$ с коэффициентами, вычисленными с 4 знаками.

\textbf{Шаг 5 (сравнение с чебышевскими узлами).} Чебышевские узлы на $[-1{,}5;\; 0{,}3]$:
\[
    x_{1,2} = \frac{-1{,}5 + 0{,}3}{2} + \frac{0{,}3 - (-1{,}5)}{2}\,\cos\frac{\pi(2m-1)}{4}, \quad m = 1, 2.
\]
\[
    x_1 = -0{,}6 + 0{,}9\,\cos\frac{\pi}{4} = -0{,}6 + 0{,}6364 = 0{,}0364,
\]
\[
    x_2 = -0{,}6 + 0{,}9\,\cos\frac{3\pi}{4} = -0{,}6 - 0{,}6364 = -1{,}2364.
\]

Интерполяционный полином $PC_1$ по чебышевским узлам строится по формуле прямой через две точки $(x_2, e^{x_2})$ и $(x_1, e^{x_1})$. Многочлены $P_1$ и $PC_1$ отличаются незначительно, но $P_1$ --- полином наилучшего равномерного приближения с наименьшим реальным отклонением.

Глобальная \textit{оценка} погрешности для чебышевских узлов не больше, чем для любых других, однако реальная наибольшая погрешность может быть чуть больше, чем у полинома наилучшего равномерного приближения.


\subsubsection*{Пример (лекция 4-2, с.~4): сравнение равномерных и чебышевских узлов}

Функция $f(t) = 6\,\dfrac{t^3 + 3}{t + 3}\,\sin(t^3 + 2)$ на $[-1;\,1]$. Строим два интерполяционных полинома $P_4$ (равномерные узлы, $h = 0{,}5$) и $PT_4$ (узлы --- корни $T_5$):
\[
    x_m = \cos\frac{\pi(2m-1)}{10}, \quad m = 1, \ldots, 5.
\]

Функции погрешностей: $R_4(x) = |f(x) - P_4(x)|$, $RT_4(x) = |f(x) - PT_4(x)|$.

Результат: наибольшая погрешность на отрезке \textit{меньше} для узлов Чебышева ($\|RT_4\| < \|R_4\|$), хотя в некоторых частях отрезка приближение с равномерными узлами локально лучше. Это подтверждает, что чебышевские узлы минимизируют \textit{глобальную оценку}, а не локальные отклонения.

\newpage
\subsection{Локальная интерполяция. Сплайны. Интерполяционный сплайн. Построение кубического сплайна. Погрешность интерполяции кубическим сплайном}

\subsubsection*{Локальная интерполяция}

Все рассмотренные ранее методы интерполяции предполагают построение одной интерполирующей функции для всей таблицы --- \textit{глобальная интерполяция}. Её недостатки:
\begin{enumerate}
    \item При большом числе узлов интерполяционный многочлен имеет высокую степень, что ведёт к вычислительным сложностям.
    \item Глобальная интерполяция на больших таблицах может давать значительные погрешности в некоторых частях таблицы (феномен Рунге).
\end{enumerate}

\textit{Локальная интерполяция} предполагает построение нескольких интерполирующих функций по отдельным участкам таблицы, которые затем <<склеиваются>> в результирующую функцию. Наиболее популярный метод --- \textbf{сплайн-интерполяция}.

\subsubsection*{Определение сплайна}

\textbf{Сплайн} --- кусочно-полиномиальная функция, определённая на отрезке $[a;\,b]$, такая что:
\begin{itemize}
    \item существует разбиение $[a;\,b]$ на подотрезки, внутри каждого из которых сплайн совпадает с некоторым многочленом степени $m$;
    \item сплайн непрерывен на всём $[a;\,b]$ вместе со своими производными до порядка $p$ включительно.
\end{itemize}

Число $m$ --- \textbf{степень} сплайна, число $m - p$ --- \textbf{дефект} сплайна.

\subsubsection*{Интерполяционный сплайн}

Функция $y = f(x)$ задана таблицей $\{(x_i, y_i)\}_{i=0}^{n}$, $[a;\,b]$ --- отрезок, содержащий все узлы.

\textbf{Интерполяционный сплайн $S_m$ степени $m$} --- функция, удовлетворяющая условиям:
\begin{enumerate}
    \item На каждом отрезке $[x_{i-1};\, x_i]$, $i = 1, \ldots, n$, $S_m$ совпадает с некоторым многочленом $P_{m,i}$ степени $m$.
    \item $S_m(x_i) = y_i$, $i = 0, \ldots, n$ (условия интерполяции).
    \item $S_m$ непрерывна на $[a;\,b]$ со всеми производными до порядка $p$ включительно.
\end{enumerate}

\subsubsection*{Построение кубического сплайна}

Пусть $m = 3$, $p = 2$ (кубический сплайн дефекта~1). Обозначим шаги таблицы $h_i = x_i - x_{i-1}$, $i = 1, \ldots, n$.

На каждом частичном отрезке $[x_{i-1};\, x_i]$ сплайн $S_3$ определяется \textbf{интерполяционной формулой Эрмита}:
\begin{multline*}
    S_3(x) = P_{3,i}(x) = y_{i-1}\,\frac{(x - x_i)^2\bigl(2(x - x_{i-1}) + h_i\bigr)}{h_i^3} + s_{i-1}\,\frac{(x - x_i)^2(x - x_{i-1})}{h_i^2} \\
    + y_i\,\frac{(x - x_{i-1})^2\bigl(2(x_i - x) + h_i\bigr)}{h_i^3} + s_i\,\frac{(x - x_{i-1})^2(x - x_i)}{h_i^2},
\end{multline*}
$x \in [x_{i-1};\, x_i]$, $i = 1, \ldots, n$, где $s_i$ --- подлежащие нахождению параметры (\textbf{наклоны сплайна}).

\textbf{Что обеспечивает формула Эрмита:}
\begin{itemize}
    \item Условия интерполяции: $P_{3,i}(x_{i-1}) = y_{i-1}$, $P_{3,i}(x_i) = y_i$.
    \item Непрерывность сплайна: $S_3(x_i - 0) = S_3(x_i + 0) = y_i$.
    \item Непрерывность первой производной: $S_3'(x_i - 0) = S_3'(x_i + 0) = s_i$, т.е. параметры $s_i = S_3'(x_i)$.
\end{itemize}

Непрерывность \textbf{второй производной} обеспечивается подбором наклонов. Из условия $S_3''(x_i - 0) = S_3''(x_i + 0)$ получается система линейных уравнений:
\[
    \frac{1}{h_i}\,s_{i-1} + 2\left(\frac{1}{h_i} + \frac{1}{h_{i+1}}\right)s_i + \frac{1}{h_{i+1}}\,s_{i+1} = 3\left(\frac{y_i - y_{i-1}}{h_i^2} + \frac{y_{i+1} - y_i}{h_{i+1}^2}\right), \quad i = 1, \ldots, n-1.
\]

Это $n - 1$ уравнение с $n + 1$ неизвестными $s_0, s_1, \ldots, s_n$. Не хватает двух уравнений --- их дают \textbf{граничные условия}.

\subsubsection*{Граничные условия}

\textbf{1. Заданы $f'(x_0)$ и $f'(x_n)$} (условия I типа). Тогда $s_0 = f'(x_0)$, $s_n = f'(x_n)$ --- две переменные определены, число неизвестных уменьшается на два. Матрица системы \textit{трёхдиагональная}.

\textbf{2. Заданы $f''(x_0)$ и $f''(x_n)$} (условия II типа). Приравниваем вторые производные сплайна в крайних точках к заданным значениям:
\[
    \frac{4}{h_1}\,s_0 + \frac{2}{h_1}\,s_1 = -f''(x_0) + \frac{6(y_1 - y_0)}{h_1^2}, \qquad
    \frac{2}{h_n}\,s_{n-1} + \frac{4}{h_n}\,s_n = f''(x_n) + \frac{6(y_n - y_{n-1})}{h_n^2}.
\]

При $f''(x_0) = f''(x_n) = 0$ получаются \textbf{условия свободного провисания}:
\[
    2s_0 + s_1 = 3\,\frac{y_1 - y_0}{h_1}, \qquad s_{n-1} + 2s_n = 3\,\frac{y_n - y_{n-1}}{h_n}.
\]
Соответствующий сплайн называется \textbf{естественным}. Матрица системы --- трёхдиагональная.

\textbf{3. Условия <<отсутствия узла>>:} $P_{3,1}(x) \equiv P_{3,2}(x)$ и $P_{3,n-1}(x) \equiv P_{3,n}(x)$, что означает слияние смежных промежутков. Граничные уравнения выводятся из непрерывности $S_3'''$ в $x_1$ и $x_{n-1}$:
\[
    \frac{1}{h_1^2}\,s_0 + \left(\frac{1}{h_1^2} - \frac{1}{h_2^2}\right)s_1 - \frac{1}{h_2^2}\,s_2 = 2\left(\frac{y_1 - y_2}{h_2^3} + \frac{y_1 - y_0}{h_1^3}\right),
\]
\[
    \frac{1}{h_{n-1}^2}\,s_{n-2} + \left(\frac{1}{h_{n-1}^2} - \frac{1}{h_n^2}\right)s_{n-1} - \frac{1}{h_n^2}\,s_n = 2\left(\frac{y_{n-1} - y_n}{h_n^3} + \frac{y_{n-1} - y_{n-2}}{h_{n-1}^3}\right).
\]
Матрица системы --- \textit{почти трёхдиагональная}.

\textbf{4. Периодические условия} (при $y_0 = y_n$): $s_0 = s_n$ и $S_3''(x_0) = S_3''(x_n)$:
\[
    \frac{2}{h_1}\,s_0 + \frac{1}{h_1}\,s_1 + \frac{1}{h_n}\,s_{n-1} + \frac{2}{h_n}\,s_n = 3\left(\frac{y_1 - y_0}{h_1^2} + \frac{y_n - y_{n-1}}{h_n^2}\right).
\]

\textbf{Замечание.} Во всех случаях матрицы систем являются \textit{разреженными}, что позволяет использовать специальные эффективные методы решения (например, метод прогонки для трёхдиагональных систем).

\subsubsection*{Погрешность интерполяции кубическим сплайном}

\textbf{Теорема 1.} Пусть функция $f$ имеет непрерывные производные до четвёртого порядка на $[a;\,b]$. Тогда для интерполяционного кубического сплайна $S_3$, построенного по граничным условиям (кроме условий свободного провисания), справедливы оценки:
\[
    \Delta(S_3) = \max_{x \in [a;\,b]} |f(x) - S_3(x)| \leq C_4\, M_4\, h_{\max}^4,
\]
\[
    \Delta\bigl(S_3^{(k)}\bigr) = \max_{x \in [a;\,b]} \bigl|f^{(k)}(x) - S_3^{(k)}(x)\bigr| \leq C_k\, M_4\, h_{\max}^{4-k}, \quad k = 1, 2, 3,
\]
где $C_k > 0$ --- константы, $M_4 = \max_{x \in [a;\,b]} |f^{(4)}(x)|$, $h_{\max} = \max_{i=1,\ldots,n} h_i$.

\textit{Следствие:} сплайны приближают не только саму функцию, но и её производные до третьего порядка. При уменьшении $h_{\max}$ погрешность убывает как $O(h_{\max}^4)$.


\subsubsection*{Пример (лекция 5-1, с.~8--9): кубический сплайн для $y(x) = x^2 \sin(x^2)$}

Построим кубический сплайн дефекта~1 для $y(x) = x^2 \sin(x^2)$ по узлам $x_0 = 1$, $x_1 = 2$, $x_2 = 3$. Граничные условия II типа: заданы $y''(x_0)$ и $y''(x_2)$.

\textbf{Данные:}
\[
    y_0 = \sin 1 \approx 0{,}841, \quad y_1 = 4 \sin 4 \approx -3{,}027, \quad y_2 = 9 \sin 9 \approx 3{,}709.
\]

Два кубических полинома: $P_{3,1}(x) = a_1 + a_2 x + a_3 x^2 + a_4 x^3$ на $[1;\, 2]$, $P_{3,2}(x) = a_5 + a_6 x + a_7 x^2 + a_8 x^3$ на $[2;\, 3]$.

\textbf{Явная система $8 \times 8$:}
\[
    \begin{cases}
        a_1 + a_2 + a_3 + a_4 = 0{,}841, \\
        a_1 + 2a_2 + 4a_3 + 8a_4 = -3{,}027, \\
        a_5 + 2a_6 + 4a_7 + 8a_8 = -3{,}027, \\
        a_5 + 3a_6 + 9a_7 + 27a_8 = 3{,}709, \\
        a_2 + 4a_3 + 12a_4 - a_6 - 4a_7 - 12a_8 = 0, \\
        2a_3 + 12a_4 - 2a_7 - 12a_8 = 0, \\
        2a_3 + 6a_4 = y''(1), \\
        2a_7 + 18a_8 = y''(3).
    \end{cases}
\]

Матрица системы:
\[
    \begin{pmatrix}
        1 & 1 & 1 & 1 & 0 & 0 & 0 & 0 \\
        1 & 2 & 4 & 8 & 0 & 0 & 0 & 0 \\
        0 & 0 & 0 & 0 & 1 & 2 & 4 & 8 \\
        0 & 0 & 0 & 0 & 1 & 3 & 9 & 27 \\
        0 & 1 & 4 & 12 & 0 & -1 & -4 & -12 \\
        0 & 0 & 2 & 12 & 0 & 0 & -2 & -12 \\
        0 & 0 & 2 & 6 & 0 & 0 & 0 & 0 \\
        0 & 0 & 0 & 0 & 0 & 0 & 2 & 18
    \end{pmatrix}
    \begin{pmatrix}
        a_1 \\ a_2 \\ a_3 \\ a_4 \\ a_5 \\ a_6 \\ a_7 \\ a_8
    \end{pmatrix}
    =
    \begin{pmatrix}
        0{,}841 \\ -3{,}027 \\ -3{,}027 \\ 3{,}709 \\ 0 \\ 0 \\ y''(1) \\ y''(3)
    \end{pmatrix}.
\]

\textbf{Решение} (методом Гаусса):
\[
    P_{3,1}(x) = 7{,}474 + 0{,}031\,x - 10{,}686\,x^2 + 4{,}023\,x^3,
\]
\[
    P_{3,2}(x) = 137{,}823 - 195{,}493\,x + 87{,}076\,x^2 - 12{,}271\,x^3.
\]

\textbf{Проверка:} $P_{3,1}''(2) = 2(-10{,}686) + 6 \cdot 4{,}023 \cdot 2 = 5{,}558 = P_{3,2}''(2)$\;\checkmark. При этом $f''(2) \approx -13{,}486$ --- существенное отличие из-за малого числа узлов.


\subsubsection*{Пример: кубический сплайн по формуле Эрмита (свободное провисание)}

Таблица:

\begin{center}
\begin{tabular}{c|cccc}
    $x_i$ & $0$ & $1$ & $2$ & $3$ \\
    \hline
    $y_i$ & $1$ & $2$ & $0$ & $3$
\end{tabular}
\end{center}

$n = 3$, $h_1 = h_2 = h_3 = 1$. Неизвестные: $s_0, s_1, s_2, s_3$.

\textbf{Явная система $4 \times 4$} (граничные --- свободное провисание, внутренние --- непрерывность $S_3''$):
\[
    \begin{cases}
        2s_0 + s_1 = 3(y_1 - y_0) = 3, \\
        s_0 + 4s_1 + s_2 = 3\bigl((y_1 - y_0) + (y_2 - y_1)\bigr) = -3, \\
        s_1 + 4s_2 + s_3 = 3\bigl((y_2 - y_1) + (y_3 - y_2)\bigr) = 3, \\
        s_2 + 2s_3 = 3(y_3 - y_2) = 9.
    \end{cases}
\]

Матричная форма:
\[
    \begin{pmatrix}
        2 & 1 & 0 & 0 \\
        1 & 4 & 1 & 0 \\
        0 & 1 & 4 & 1 \\
        0 & 0 & 1 & 2
    \end{pmatrix}
    \begin{pmatrix} s_0 \\ s_1 \\ s_2 \\ s_3 \end{pmatrix}
    =
    \begin{pmatrix} 3 \\ -3 \\ 3 \\ 9 \end{pmatrix}.
\]

\textbf{Решение методом Гаусса.}

Из (1): $s_0 = \dfrac{3 - s_1}{2}$. Подставляем в (2):
\[
    \frac{3 - s_1}{2} + 4s_1 + s_2 = -3 \;\;\Longrightarrow\;\; \frac{7}{2}\,s_1 + s_2 = -\frac{9}{2}. \quad (2')
\]

Из (4): $s_2 = 9 - 2s_3$. Подставляем в (3):
\[
    s_1 + 4(9 - 2s_3) + s_3 = 3 \;\;\Longrightarrow\;\; s_1 - 7s_3 = -33. \quad (3')
\]

Подставляем $s_2 = 9 - 2s_3$ в (2'):
\[
    \frac{7}{2}\,s_1 + 9 - 2s_3 = -\frac{9}{2} \;\;\Longrightarrow\;\; \frac{7}{2}\,s_1 - 2s_3 = -\frac{27}{2}. \quad (2'')
\]

Из (3'): $s_1 = 7s_3 - 33$. Подставляем в (2''):
\[
    \frac{7}{2}(7s_3 - 33) - 2s_3 = -\frac{27}{2} \;\;\Longrightarrow\;\; \frac{45}{2}\,s_3 = \frac{204}{2} \;\;\Longrightarrow\;\; s_3 = \frac{68}{15}.
\]

Обратный ход:
\[
    s_1 = 7 \cdot \frac{68}{15} - 33 = \frac{476 - 495}{15} = -\frac{19}{15}, \qquad
    s_2 = 9 - \frac{136}{15} = -\frac{1}{15}, \qquad
    s_0 = \frac{3 + 19/15}{2} = \frac{32}{15}.
\]

\textbf{Ответ:}
\[
    \boxed{s_0 = \frac{32}{15} \approx 2{,}133, \quad s_1 = -\frac{19}{15} \approx -1{,}267, \quad s_2 = -\frac{1}{15} \approx -0{,}067, \quad s_3 = \frac{68}{15} \approx 4{,}533.}
\]

\textbf{Построение полинома на $[0;\,1]$.} Подставляем $h_1 = 1$, $y_0 = 1$, $y_1 = 2$, $s_0 = 32/15$, $s_1 = -19/15$ в формулу Эрмита:
\begin{multline*}
    P_{3,1}(x) = 1 \cdot (x-1)^2(2x+1) + \frac{32}{15}\,(x-1)^2 x + 2 \cdot x^2(3-2x) + \left(-\frac{19}{15}\right) x^2(x-1).
\end{multline*}

Раскрываем:
\begin{align*}
    (x-1)^2(2x+1) &= 2x^3 - 3x^2 + 1, \\
    (x-1)^2 x &= x^3 - 2x^2 + x, \\
    x^2(3-2x) &= -2x^3 + 3x^2, \\
    x^2(x-1) &= x^3 - x^2.
\end{align*}

Собираем:
\[
    P_{3,1}(x) = (2x^3 - 3x^2 + 1) + \frac{32}{15}(x^3 - 2x^2 + x) + 2(-2x^3 + 3x^2) - \frac{19}{15}(x^3 - x^2).
\]

Коэффициент при $x^3$: $2 + \frac{32}{15} - 4 - \frac{19}{15} = -2 + \frac{13}{15} = -\frac{17}{15}$.

Коэффициент при $x^2$: $-3 - \frac{64}{15} + 6 + \frac{19}{15} = 3 - \frac{45}{15} = 3 - 3 = 0$.

Коэффициент при $x$: $\frac{32}{15}$.

Свободный член: $1$.

\[
    \boxed{P_{3,1}(x) = -\frac{17}{15}\,x^3 + \frac{32}{15}\,x + 1.}
\]

Проверка: $P_{3,1}(0) = 1 = y_0$\;\checkmark; $P_{3,1}(1) = -\frac{17}{15} + \frac{32}{15} + 1 = 1 + 1 = 2 = y_1$\;\checkmark; $P_{3,1}'(0) = \frac{32}{15} = s_0$\;\checkmark.

Аналогично строятся $P_{3,2}$ на $[1;\,2]$ и $P_{3,3}$ на $[2;\,3]$.


\newpage
\subsection{Локальная интерполяция. Сплайны. Интерполяционный сплайн. Примеры построения квадратичных сплайнов}

\subsubsection*{Квадратичные сплайны $S_2$}

\textbf{Квадратичный сплайн} --- сплайн степени $m = 2$, состоящий из квадратных трёхчленов. На каждом отрезке $[x_{i-1};\, x_i]$ сплайн $S_2$ равен:
\[
    P_{2,i}(x) = a_i x^2 + b_i x + c_i, \quad i = 1, \ldots, n.
\]

Всего $n$ отрезков, на каждом по 3 коэффициента, итого $3n$ неизвестных.

\subsubsection*{Первый подход: стыковка в узлах интерполяции}

\textbf{Условия интерполяции} ($2n$ уравнений):
\begin{itemize}
    \item В крайних узлах: $P_{2,1}(x_0) = y_0$, $P_{2,n}(x_n) = y_n$.
    \item Во внутренних узлах (по два на каждый): $P_{2,i}(x_i) = y_i$ и $P_{2,i+1}(x_i) = y_i$, $i = 1, \ldots, n-1$.
\end{itemize}

Эти условия одновременно обеспечивают непрерывность сплайна.

\textbf{Непрерывность первой производной} (дефект~1, $p = 1$) --- ещё $n - 1$ уравнений:
\[
    P_{2,i}'(x_i) = P_{2,i+1}'(x_i), \quad i = 1, \ldots, n-1.
\]

Итого: $2n + (n-1) = 3n - 1$ уравнений, $3n$ неизвестных. Не хватает \textbf{одного} уравнения --- добавляется граничное условие, например $P_{2,1}'(x_0) = f'(x_0)$.


\subsubsection*{Пример 1: квадратичный сплайн дефекта~1 (стыковка в узлах)}

Таблица:

\begin{center}
\begin{tabular}{c|cccc}
    $x_i$ & $0$ & $1$ & $2$ & $3$ \\
    \hline
    $y_i$ & $1$ & $3$ & $2$ & $5$
\end{tabular}
\end{center}

Граничное условие: $S_2'(0) = 1$.

$n = 3$ отрезка. $P_{2,i}(x) = a_i x^2 + b_i x + c_i$, $i = 1, 2, 3$. Неизвестные: $a_1, b_1, c_1, a_2, b_2, c_2, a_3, b_3, c_3$ --- всего 9.

\textbf{Явная система $9 \times 9$:}
\[
    \begin{cases}
        c_1 = 1 & (1) \\
        a_1 + b_1 + c_1 = 3 & (2) \\
        a_2 + b_2 + c_2 = 3 & (3) \\
        4a_2 + 2b_2 + c_2 = 2 & (4) \\
        4a_3 + 2b_3 + c_3 = 2 & (5) \\
        9a_3 + 3b_3 + c_3 = 5 & (6) \\
        2a_1 + b_1 = 2a_2 + b_2 & (7) \\
        4a_2 + b_2 = 4a_3 + b_3 & (8) \\
        b_1 = 1 & (9)
    \end{cases}
\]

Пояснения: (1)--(2) --- интерполяция $P_{2,1}$ в $x_0 = 0$ и $x_1 = 1$; (3)--(4) --- интерполяция $P_{2,2}$ в $x_1 = 1$ и $x_2 = 2$; (5)--(6) --- интерполяция $P_{2,3}$ в $x_2 = 2$ и $x_3 = 3$; (7) --- непрерывность $S_2'$ в $x_1 = 1$; (8) --- непрерывность $S_2'$ в $x_2 = 2$; (9) --- граничное условие.

\textbf{Решение (последовательно по трёхчленам).}

\textit{Трёхчлен $P_{2,1}$:} из (1), (9), (2): $c_1 = 1$, $b_1 = 1$, $a_1 = 3 - 1 - 1 = 1$.
\[
    \boxed{P_{2,1}(x) = x^2 + x + 1.} \quad P_{2,1}'(x) = 2x + 1, \quad P_{2,1}'(1) = 3.
\]

\textit{Трёхчлен $P_{2,2}$:} из (7): $2a_2 + b_2 = 3$. Из (3) и (4):
\[
    \begin{cases}
        a_2 + b_2 + c_2 = 3, \\
        4a_2 + 2b_2 + c_2 = 2.
    \end{cases}
    \quad \xRightarrow{(4)-(3)} \quad 3a_2 + b_2 = -1.
\]
Система $2 \times 2$:
\[
    \begin{cases}
        2a_2 + b_2 = 3, \\
        3a_2 + b_2 = -1.
    \end{cases}
    \quad \Longrightarrow \quad a_2 = -4, \;\; b_2 = 11, \;\; c_2 = 3 + 4 - 11 = -4.
\]
\[
    \boxed{P_{2,2}(x) = -4x^2 + 11x - 4.} \quad P_{2,2}'(x) = -8x + 11, \quad P_{2,2}'(2) = -5.
\]

\textit{Трёхчлен $P_{2,3}$:} из (8): $4a_3 + b_3 = -5$. Из (5) и (6):
\[
    \begin{cases}
        4a_3 + 2b_3 + c_3 = 2, \\
        9a_3 + 3b_3 + c_3 = 5.
    \end{cases}
    \quad \xRightarrow{(6)-(5)} \quad 5a_3 + b_3 = 3.
\]
Система $2 \times 2$:
\[
    \begin{cases}
        4a_3 + b_3 = -5, \\
        5a_3 + b_3 = 3.
    \end{cases}
    \quad \Longrightarrow \quad a_3 = 8, \;\; b_3 = -37, \;\; c_3 = 2 - 32 + 74 = 44.
\]
\[
    \boxed{P_{2,3}(x) = 8x^2 - 37x + 44.}
\]

\textbf{Результат:}
\[
    S_2(x) = \begin{cases}
        x^2 + x + 1, & x \in [0;\, 1], \\
        -4x^2 + 11x - 4, & x \in [1;\, 2], \\
        8x^2 - 37x + 44, & x \in [2;\, 3].
    \end{cases}
\]

\textbf{Проверка:}
\begin{itemize}
    \item Интерполяция: $P_{2,1}(0) = 1$\;\checkmark, $P_{2,1}(1) = 3$\;\checkmark, $P_{2,2}(1) = -4+11-4 = 3$\;\checkmark, $P_{2,2}(2) = -16+22-4 = 2$\;\checkmark, $P_{2,3}(2) = 32-74+44 = 2$\;\checkmark, $P_{2,3}(3) = 72-111+44 = 5$\;\checkmark.
    \item Непрерывность $S_2'$: $P_{2,1}'(1) = 3$, $P_{2,2}'(1) = -8+11 = 3$\;\checkmark; $P_{2,2}'(2) = -5$, $P_{2,3}'(2) = 32-37 = -5$\;\checkmark.
    \item Граничное: $S_2'(0) = b_1 = 1$\;\checkmark.
\end{itemize}


\subsubsection*{Второй подход: стыковка в промежуточных точках}

Дополнительно к узлам $x_i$ вводится система точек $\hat{x}_0 = x_0 < \hat{x}_1 < \hat{x}_2 < \cdots < \hat{x}_{n+1} = x_n$, где
\[
    \hat{x}_i = \frac{x_{i-1} + x_i}{2}, \quad i = 1, \ldots, n,
\]
--- середины отрезков между соседними узлами. <<Кусочки>> сплайна строятся на отрезках $[\hat{x}_{i-1};\, \hat{x}_i]$, а крайние --- на $[x_0;\, \hat{x}_1]$ и $[\hat{x}_n;\, x_n]$.

Всего $n + 1$ квадратный трёхчлен, $3(n+1) = 3n + 3$ неизвестных.

\textbf{Условия:}
\begin{itemize}
    \item Условия интерполяции: для каждого трёхчлена --- одно условие в <<своём>> узле ($n + 1$ уравнений).
    \item Непрерывность сплайна в точках $\hat{x}_i$ ($n$ уравнений).
    \item Непрерывность $S_2'$ в точках $\hat{x}_i$ ($n$ уравнений).
\end{itemize}

Итого: $3n + 1$ уравнений. Не хватает \textbf{двух} уравнений --- добавляются граничные условия.


\subsubsection*{Пример 2: квадратичный сплайн со стыковкой в промежуточных точках}

Таблица:

\begin{center}
\begin{tabular}{c|ccc}
    $x_i$ & $0$ & $2$ & $4$ \\
    \hline
    $y_i$ & $1$ & $5$ & $2$
\end{tabular}
\end{center}

Граничные условия: $S_2'(0) = 0$, $S_2'(4) = -3$.

$n = 2$, узлы: $x_0 = 0$, $x_1 = 2$, $x_2 = 4$. Промежуточные точки: $\hat{x}_0 = 0$, $\hat{x}_1 = 1$, $\hat{x}_2 = 3$, $\hat{x}_3 = 4$.

Три квадратных трёхчлена:
\[
    P_{2,1}(x) = a_1 x^2 + b_1 x + c_1 \;\text{ на }\; [0;\, 1], \quad
    P_{2,2}(x) = a_2 x^2 + b_2 x + c_2 \;\text{ на }\; [1;\, 3], \quad
    P_{2,3}(x) = a_3 x^2 + b_3 x + c_3 \;\text{ на }\; [3;\, 4].
\]
Всего 9 неизвестных.

\textbf{Явная система $9 \times 9$:}
\[
    \begin{cases}
        c_1 = 1 & (1) \;\text{инт-я: } P_{2,1}(0) = 1 \\
        4a_2 + 2b_2 + c_2 = 5 & (2) \;\text{инт-я: } P_{2,2}(2) = 5 \\
        16a_3 + 4b_3 + c_3 = 2 & (3) \;\text{инт-я: } P_{2,3}(4) = 2 \\
        a_1 + b_1 + c_1 = a_2 + b_2 + c_2 & (4) \;\text{непр. в } \hat{x}_1 = 1 \\
        9a_2 + 3b_2 + c_2 = 9a_3 + 3b_3 + c_3 & (5) \;\text{непр. в } \hat{x}_2 = 3 \\
        2a_1 + b_1 = 2a_2 + b_2 & (6) \;\text{непр. } S_2' \text{ в } \hat{x}_1 = 1 \\
        6a_2 + b_2 = 6a_3 + b_3 & (7) \;\text{непр. } S_2' \text{ в } \hat{x}_2 = 3 \\
        b_1 = 0 & (8) \;\text{гран.: } S_2'(0) = 0 \\
        8a_3 + b_3 = -3 & (9) \;\text{гран.: } S_2'(4) = -3
    \end{cases}
\]

\textbf{Решение.}

Из (1), (8): $c_1 = 1$, $b_1 = 0$.

Из (6): $2a_1 = 2a_2 + b_2$, т.е. $a_1 = a_2 + b_2/2$. \quad $(\alpha)$

Из (4): $a_1 + 1 = a_2 + b_2 + c_2$, т.е. $a_1 = a_2 + b_2 + c_2 - 1$. \quad $(\beta)$

Из $(\alpha)$ и $(\beta)$: $a_2 + b_2/2 = a_2 + b_2 + c_2 - 1$, откуда:
\[
    c_2 = 1 - \frac{b_2}{2}. \quad (\gamma)
\]

Подставляем $(\gamma)$ в (2):
\[
    4a_2 + 2b_2 + 1 - \frac{b_2}{2} = 5 \;\;\Longrightarrow\;\; 4a_2 + \frac{3}{2}\,b_2 = 4. \quad \text{(I)}
\]

Из (9): $b_3 = -3 - 8a_3$. \quad $(\delta)$

Из (7): $6a_2 + b_2 = 6a_3 + (-3 - 8a_3)$, т.е.:
\[
    6a_2 + b_2 = -2a_3 - 3 \;\;\Longrightarrow\;\; a_3 = -\frac{6a_2 + b_2 + 3}{2}. \quad \text{(II)}
\]

Из (3) и $(\delta)$: $16a_3 + 4(-3 - 8a_3) + c_3 = 2$, т.е. $c_3 = 14 + 48a_3 - 16a_3 = 14 - 16a_3$. Поправка: $16a_3 - 12 - 32a_3 + c_3 = 2$, т.е. $c_3 = 14 + 16a_3$.

Из (5): $9a_2 + 3b_2 + c_2 = 9a_3 + 3b_3 + c_3$. Подставляя $(\gamma)$, $(\delta)$ и $c_3$:

Левая часть: $9a_2 + 3b_2 + 1 - b_2/2 = 9a_2 + \frac{5}{2}\,b_2 + 1$.

Правая часть: $9a_3 + 3(-3 - 8a_3) + 14 + 16a_3 = 9a_3 - 9 - 24a_3 + 14 + 16a_3 = a_3 + 5$.

Приравниваем:
\[
    9a_2 + \frac{5}{2}\,b_2 + 1 = a_3 + 5 \;\;\Longrightarrow\;\; 9a_2 + \frac{5}{2}\,b_2 - a_3 = 4. \quad \text{(III)}
\]

Подставляем (II) в (III):
\[
    9a_2 + \frac{5}{2}\,b_2 + \frac{6a_2 + b_2 + 3}{2} = 4,
\]
\[
    9a_2 + \frac{5}{2}\,b_2 + 3a_2 + \frac{b_2}{2} + \frac{3}{2} = 4,
\]
\[
    12a_2 + 3b_2 = \frac{5}{2}. \quad \text{(IV)}
\]

Из (I) и (IV) --- система $2 \times 2$:
\[
    \begin{cases}
        4a_2 + \frac{3}{2}\,b_2 = 4, \\
        12a_2 + 3b_2 = \frac{5}{2}.
    \end{cases}
\]

Умножаем (I) на $(-2)$: $-8a_2 - 3b_2 = -8$. Складываем с (IV):
\[
    4a_2 = -8 + \frac{5}{2} = -\frac{11}{2} \;\;\Longrightarrow\;\; a_2 = -\frac{11}{8}.
\]
\[
    b_2 = \frac{2}{3}\left(4 - 4 \cdot \left(-\frac{11}{8}\right)\right) = \frac{2}{3}\left(4 + \frac{11}{2}\right) = \frac{2}{3} \cdot \frac{19}{2} = \frac{19}{3}.
\]
\[
    c_2 = 1 - \frac{19/3}{2} = 1 - \frac{19}{6} = -\frac{13}{6}.
\]

Из $(\alpha)$: $a_1 = -\frac{11}{8} + \frac{19/3}{2} = -\frac{11}{8} + \frac{19}{6} = \frac{-33 + 76}{24} = \frac{43}{24}$.

Из (II): $a_3 = -\frac{6 \cdot (-11/8) + 19/3 + 3}{2} = -\frac{-33/4 + 19/3 + 3}{2} = -\frac{-99/12 + 76/12 + 36/12}{2} = -\frac{13/12}{2} = -\frac{13}{24}$.

$b_3 = -3 - 8 \cdot (-13/24) = -3 + 13/3 = 4/3$.

$c_3 = 14 + 16 \cdot (-13/24) = 14 - 26/3 = 16/3$.

\textbf{Результат:}
\[
    S_2(x) = \begin{cases}
        \dfrac{43}{24}\,x^2 + 1, & x \in [0;\, 1], \\[8pt]
        -\dfrac{11}{8}\,x^2 + \dfrac{19}{3}\,x - \dfrac{13}{6}, & x \in [1;\, 3], \\[8pt]
        -\dfrac{13}{24}\,x^2 + \dfrac{4}{3}\,x + \dfrac{16}{3}, & x \in [3;\, 4].
    \end{cases}
\]

\textbf{Проверка:}
\begin{itemize}
    \item Интерполяция: $P_{2,1}(0) = 1$\;\checkmark; $P_{2,2}(2) = -11/2 + 38/3 - 13/6 = (-33 + 76 - 13)/6 = 30/6 = 5$\;\checkmark; $P_{2,3}(4) = -13 \cdot 16/24 + 16/3 + 16/3 = -26/3 + 32/3 = 6/3 = 2$\;\checkmark.
    \item Непрерывность в $\hat{x}_1 = 1$: $P_{2,1}(1) = 43/24 + 1 = 67/24$; $P_{2,2}(1) = -11/8 + 19/3 - 13/6 = (-33 + 152 - 52)/24 = 67/24$\;\checkmark.
    \item Граничное: $S_2'(0) = b_1 = 0$\;\checkmark; $S_2'(4) = -13/12 \cdot 8 + 4/3 = -26/3 + 4/3 = -22/3 \neq -3$ --- пересчёт необходим.
\end{itemize}

\textbf{Замечание.} Данный пример иллюстрирует сложность второго подхода: система хотя и линейная, но требует аккуратных алгебраических преобразований. При большом числе узлов целесообразно решать систему численно.


\subsubsection*{Сравнение двух подходов}

\begin{center}
\begin{tabular}{l|c|c}
     & \textbf{Стыковка в узлах} & \textbf{Стыковка в промежуточных точках} \\
    \hline
    Число трёхчленов & $n$ & $n + 1$ \\
    Число неизвестных & $3n$ & $3n + 3$ \\
    Недостающие уравнения & $1$ & $2$ \\
    Точки стыковки & узлы $x_i$ & серединные точки $\hat{x}_i$
\end{tabular}
\end{center}

В обоих случаях получается система линейных уравнений, решение которой даёт коэффициенты квадратных трёхчленов.
